\documentclass[11pt,a4paper]{report}
% Compilation : tectonic Cours_Contraintes.tex
% Figures à données : python3 tools/make_figures.py (lit Code/resultats/contraintes_*.json,
% produits par les scripts de Code/contraintes/, voir Code/contraintes/README.md)
% Préambule commun (repris du cours IA et QUBO)
\usepackage{iftex}
\ifPDFTeX
  \usepackage[utf8]{inputenc}
  \usepackage[T1]{fontenc}
  \usepackage{lmodern}
\else
  \usepackage{fontspec}
\fi
\usepackage[french]{babel}
\usepackage[a4paper,margin=2.3cm]{geometry}
\usepackage{amsmath,amssymb,mathtools}
\usepackage{xcolor}
\usepackage{tikz}
\usetikzlibrary{arrows.meta,calc,positioning,shapes.geometric,decorations.pathreplacing,decorations.markings,patterns,babel}
\usepackage{pgfplots}
\pgfplotsset{compat=1.18}
\usepgfplotslibrary{groupplots}
\usepackage{tabularx,booktabs,array}
\usepackage{enumitem}
\setlist{nosep,leftmargin=*}
\usepackage{listings}
\usepackage[font=small,labelfont=bf]{caption}
\usepackage[raggedright]{titlesec}
\usepackage[most]{tcolorbox}
\usepackage[colorlinks=true,linkcolor=blue!50!black,urlcolor=blue!50!black,citecolor=blue!50!black]{hyperref}

\definecolor{cblue}{HTML}{2A78D6}
\definecolor{corange}{HTML}{EB6834}
\definecolor{cgreen}{HTML}{1BAF7A}
\definecolor{cink}{HTML}{1F2328}
\definecolor{cmuted}{HTML}{6B6F76}
\definecolor{cteal}{HTML}{0E8A8A}

% ---------------------------------------------------------------- notations
\newcommand{\dd}{\mathrm{d}}
\newcommand{\Ue}{U_e}
\renewcommand{\Re}{\mathit{Re}}
\newcommand{\vect}[1]{\mathbf{#1}}
\newcommand{\code}[1]{\texttt{#1}}

% ---------------------------------------------------------------- encadrés
\tcbset{coursbox/.style={enhanced,breakable,boxrule=0pt,leftrule=2.5pt,arc=0pt,
  left=6pt,right=6pt,top=4pt,bottom=4pt,fonttitle=\bfseries\small,coltitle=cink,
  attach title to upper={\ }}}
\newtcolorbox{definition}[1][]{coursbox,colback=cblue!5,colframe=cblue,title={Définition#1.}}
\newtcolorbox{resultat}[1][]{coursbox,colback=cgreen!6,colframe=cgreen,title={À retenir#1.}}
\newtcolorbox{projet}[1][]{coursbox,colback=corange!6,colframe=corange,title={Dans notre projet#1.}}
\newtcolorbox{piege}[1][]{coursbox,colback=cmuted!7,colframe=cmuted,title={Piège#1.}}
\newtcolorbox{codebox}[1][]{coursbox,colback=cmuted!3,colframe=cblue!60,title={Dans le code#1.}}

% ---------------------------------------------------------------- exercices
\newcounter{exo}[chapter]
\renewcommand{\theexo}{\thechapter.\arabic{exo}}
\newenvironment{exercice}[1][]{\refstepcounter{exo}\par\medskip\noindent
  \textbf{\color{cblue}Exercice~\theexo}\if\relax\detokenize{#1}\relax\else\ (#1)\fi.\ }{\par}
\newcommand{\exercices}{\section*{Exercices}\addcontentsline{toc}{section}{Exercices}
  \noindent\textit{Corrigés en annexe~\ref{annexe:corriges}. Chercher d'abord sans les regarder.}}

\tikzset{
  flow/.style={draw=cink,rounded corners=2pt,fill=white,align=center,font=\small,inner sep=4pt},
  arr/.style={-{Stealth[length=2.2mm]},cink,thick},
}

\lstset{basicstyle=\ttfamily\footnotesize,breaklines=true,columns=fullflexible,frame=single,rulecolor=\color{cmuted!50},
  backgroundcolor=\color{cmuted!6},keywordstyle=\color{cblue}\bfseries,commentstyle=\color{cmuted}\itshape,
  stringstyle=\color{corange},showstringspaces=false,language=Python,
  literate={é}{\'e}1{è}{\`e}1{à}{\`a}1{ê}{\^e}1{ç}{\c c}1{ô}{\^o}1{î}{\^i}1{ù}{\`u}1{É}{\'E}1{→}{$\to$}1}

\setcounter{secnumdepth}{2}
\setcounter{tocdepth}{1}
\setlength{\parindent}{0pt}
\setlength{\emergencystretch}{3em}
\setlength{\parskip}{0.45em}


\definecolor{cviolet}{HTML}{7A4FB3}
\newtcolorbox{pistes}[1][]{coursbox,colback=cviolet!6,colframe=cviolet,title={Pistes d'amélioration#1.}}
\newcommand{\dB}{\,\mathrm{dB}}
\newcommand{\R}{\mathbb{R}}
\newcommand{\Id}{\mathbf{I}}


% ============================================================================
% Itérations du projet HackathonQuantique : couleurs et marqueurs communs.
% Inclus par tous les cours et par la fiche projet, pour qu'une itération ait
% la même couleur partout. Pour ajouter une itération : une couleur, un nom,
% une date, une ligne dans \legendeiterations.
% ============================================================================
\definecolor{itV0}{HTML}{8D6E63}   % brun   : V0, modèle de la note d'intention
\definecolor{itV1}{HTML}{C2185B}   % framboise : V1, chaîne complète conforme à la note
\definecolor{itV2}{HTML}{00838F}   % cyan foncé : V2, problème reformulé
\definecolor{itV3}{HTML}{E09A00}   % ambre  : V3, concurrent équitable et problème plus grand
\definecolor{itV4}{HTML}{5E35B1}   % indigo : V4, réservé

\makeatletter
\newcommand{\itername}[1]{\@nameuse{itname@#1}}
\newcommand{\iterdate}[1]{\@nameuse{itdate@#1}}
\newcommand{\defiter}[3]{\@namedef{itname@#1}{#2}\@namedef{itdate@#1}{#3}}
\makeatother
\defiter{V0}{Modèle de la note d'intention}{24--30 oct. 2026}
\defiter{V1}{Chaîne complète conforme à la note}{4--5 nov. 2026}
\defiter{V2}{Problème reformulé}{5--6 nov. 2026}
\defiter{V3}{Concurrent équitable, problème plus grand}{à partir du 6 nov. 2026}

% Pastille : \iterbadge{V1}
\newcommand{\iterbadge}[1]{\tikz[baseline=(b.base)]\node[fill=it#1,text=white,rounded corners=2pt,
  inner xsep=3pt,inner ysep=1.5pt,font=\sffamily\bfseries\scriptsize] (b) {#1};}

% Bandeau de début de chapitre : \iterbanner{V0}
\newcommand{\iterbanner}[1]{\par\noindent\begin{tikzpicture}
  \node[fill=it#1!12,draw=it#1,line width=0.6pt,rounded corners=2pt,text width=\dimexpr\linewidth-12pt\relax,
    inner sep=4pt,anchor=west,font=\small] at (0,0)
    {\iterbadge{#1}\ \textbf{\color{it#1}Itération #1} --- \itername{#1} \hfill{\color{it#1}\iterdate{#1}}};
  \end{tikzpicture}\par\medskip}

% Encadré d'évolution : \begin{evolution}{V1}{titre} ... \end{evolution}
\newtcolorbox{evolution}[2]{enhanced,breakable,boxrule=0pt,leftrule=3pt,arc=0pt,
  left=6pt,right=6pt,top=4pt,bottom=4pt,colback=it#1!7,colframe=it#1,
  fonttitle=\bfseries\small,coltitle=it#1,attach title to upper={\ },
  title={\iterbadge{#1}\ Évolution #1 : #2.}}

% Légende des itérations, pour les avant-propos
\newcommand{\legendeiterations}{%
\begin{center}\small
\begin{tabular}{lll}
\iterbadge{V0} & \itername{V0} & \iterdate{V0}\\
\iterbadge{V1} & \itername{V1} & \iterdate{V1}\\
\iterbadge{V2} & \itername{V2} & \iterdate{V2}\\
\iterbadge{V3} & \itername{V3} & \iterdate{V3}\\
\end{tabular}
\end{center}}

% Encadré « pistes d'amélioration » : défini ici s'il manque (cours V0 sans cet encadré)
\providecolor{cviolet}{HTML}{7A4FB3}
\ifcsname pistes\endcsname\else
\newtcolorbox{pistes}[1][]{enhanced,breakable,boxrule=0pt,leftrule=2.5pt,arc=0pt,
  left=6pt,right=6pt,top=4pt,bottom=4pt,fonttitle=\bfseries\small,coltitle=cink,attach title to upper={\ },
  colback=cviolet!6,colframe=cviolet,title={Pistes d'amélioration#1.}}
\fi


\title{\textbf{Contraintes et pénalités}\\[0.4em]\large Comment une boucle QUBO doit traiter les contraintes :\\ pénalités, lagrangien vérifié, classifieur de convergence,\\ et l'interface des solveurs industriels (lot L7 de la V3)}
\author{Notes de travail, projet \emph{HackathonQuantique}}
\date{Version du 28 novembre 2026}

\begin{document}
\pagenumbering{Alph}
\maketitle
\pagenumbering{roman}

\chapter*{Avant-propos}
\addcontentsline{toc}{chapter}{Avant-propos}
\paragraph{Objet.} Les lots précédents ont montré, chacun à sa manière, que les contraintes coûtent plus cher à la boucle QUBO que le choix du modèle : une contrainte oubliée a détourné une campagne entière, et même corrigée, la boucle dépense plus de la moitié de son budget sur des profils infaisables. Ce cours reprend la question de zéro : quelles sont nos contraintes, comment la pratique industrielle et la recherche les traitent, et ce que cela donne dans notre boucle.

\paragraph{Ce qu'on suppose connu.} Le QUBO et ses pénalités (cours de calcul quantique, chapitres~8 et 15), la boucle et le problème CST à 30 bits (cours IA et QUBO, chapitres~6, 8 et 9). Le lagrangien, la régression logistique et le format des solveurs contraints sont posés depuis le début.

\paragraph{Fil conducteur.} Le constat (ch.~\ref{ch:constat}) $\to$ classer les contraintes avant de les traiter (ch.~\ref{ch:taxonomie}) $\to$ ce que valent les pénalités (ch.~\ref{ch:penalites}) $\to$ le lagrangien vérifié et la géométrie exactement affine (ch.~\ref{ch:lagrangien}) $\to$ le classifieur de convergence d'XFOIL (ch.~\ref{ch:classifieur}) $\to$ les résultats obtenus sans XFOIL (ch.~\ref{ch:hors-ligne}) $\to$ l'interface des solveurs industriels (ch.~\ref{ch:cqm}) $\to$ la campagne avec XFOIL (ch.~\ref{ch:campagne}) $\to$ bilan (ch.~\ref{ch:bilan}).

\paragraph{Code et chiffres.} Tout le code est dans \code{Code/contraintes/} ; la boucle de \code{Code/ia\_qubo/} reçoit seulement de nouveaux modes, l'approche actuelle reste identique (vérifié). Les figures sont produites par \code{tools/make\_figures.py} à partir des fichiers \code{Code/resultats/contraintes\_*.json}. Chaque chapitre se termine par des exercices (corrigés en annexe~\ref{annexe:corriges}) et des pistes d'amélioration.

\paragraph{Itérations du projet.} Tout ce cours appartient à l'itération V3 (lot L7).
\legendeiterations

\tableofcontents
\cleardoublepage
\pagenumbering{arabic}

\chapter{Le constat}\label{ch:constat}
\iterbanner{V3}

\section{Quatre signes, une même cause}
Chaque lot de la V3 a rencontré les contraintes, sans les traiter comme un sujet en soi :
\begin{enumerate}
\item \textbf{Le budget part en infaisables.} Dans la seconde campagne CST (cours IA et QUBO, chapitre~9), l'approche actuelle ne propose que 43~\% de profils faisables ; le reste se partage entre échecs d'XFOIL (28~\%) et $C_{L,\max}$ trop faibles (29~\%). Le hasard, lui, en propose 76~\%. Plus de la moitié des simulations sert à découvrir des contraintes.
\item \textbf{Une contrainte oubliée détourne tout.} Dans la première campagne CST, la géométrie n'était ni modélisée ni filtrée : 73~\% du budget est parti en profils qu'aucun solveur n'aurait dû recevoir.
\item \textbf{Les pénalités écrasent l'objectif.} Dans le QUBO de la V2, l'objectif n'occupe que 1,7~\% de l'étendue des énergies (cours de calcul quantique, chapitre~15) : pour un solveur quantique, le problème est presque entièrement fait de pénalités.
\item \textbf{Un modèle de contrainte trop sévère exclut le front.} Le réseau de neurones de la V2 prédisait la contrainte de manœuvre trop fidèlement, et il refusait les profils du front, qui sont au bord.
\end{enumerate}
Dans les quatre cas, on a traité des contraintes de natures très différentes avec un seul outil : la pénalité quadratique, ou rien. Ce cours sépare les contraintes selon leur nature (chapitre~\ref{ch:taxonomie}), puis donne à chacune le traitement que la pratique industrielle lui réserve.

\section{Le problème}
On garde le problème CST du lot L6 : 10 coefficients de Kulfan à 8 niveaux (30 bits), deux objectifs (traînée en croisière, SER en secteur S2), et quatre contraintes :
\begin{center}
\renewcommand{\arraystretch}{1.2}
\begin{tabularx}{\linewidth}{>{\raggedright\arraybackslash}p{4.4cm} >{\raggedright\arraybackslash}X}
\toprule
contrainte & ce qu'il faut pour la connaître \\
\midrule
faces sans croisement & la géométrie seule, quelques microsecondes \\
épaisseur maximale $\ge8~\%$ & la géométrie seule \\
$C_{L,\max}\ge1{,}2$ (manœuvre) & une simulation XFOIL en rampe d'incidence \\
convergence d'XFOIL en croisière & une simulation ; on sait seulement si elle a abouti \\
\bottomrule
\end{tabularx}
\end{center}
Le problème V2 (4 variables, 1\,024 profils, front exact connu) sert de banc d'essai : toutes ses simulations sont tabulées, et l'on peut y comparer des formulations sans lancer XFOIL.

\section{Ce que l'on compare}
\begin{itemize}
\item \textbf{Bits d'écart} (\code{modele}) : l'approche actuelle, pénalité $P(\hat C_{L,\max}-1{,}2-s)^2$ avec $s\ge0$ codé sur 3 qubits.
\item \textbf{Pénalité déséquilibrée} (\code{desequilibree}) : une pénalité sans bit d'écart proposée pour les machines quantiques (Montañez-Barrera et al., 2024).
\item \textbf{Lagrangien vérifié} (\code{lagrange}) : un terme linéaire dans le QUBO, des multiplicateurs ajustés en quelques résolutions, et une vérification classique de chaque échantillon.
\item \textbf{Lagrangien vérifié et classifieur} (\code{lagrange\_clf}) : la même chose, plus une probabilité de convergence d'XFOIL apprise.
\end{itemize}
Le modèle d'objectif est toujours le modèle quadratique en $z$, pour n'isoler que l'effet de la contrainte.

\begin{resultat}
\begin{itemize}
\item Plus de la moitié du budget de la boucle part en profils infaisables, et les pénalités occupent 98~\% du spectre du QUBO.
\item On traitait toutes les contraintes de la même façon. Ce cours les sépare selon leur nature, puis compare quatre formulations à modèle d'objectif identique.
\end{itemize}
\end{resultat}

\exercices
\begin{exercice}\label{ex:constat-budget}
Sur 300 simulations, l'approche actuelle en dépense 28~\% en échecs d'XFOIL et 29~\% en $C_{L,\max}$ trop faibles. (a) Combien de profils faisables obtient-elle ? (b) Si un meilleur traitement des contraintes divisait par deux ces deux pertes, combien de profils faisables obtiendrait-on, et à quel budget de l'approche actuelle cela correspond-il ? (c) Pourquoi ne peut-on pas espérer 100~\% de faisables, même avec des modèles parfaits ?
\end{exercice}

\begin{pistes}
\begin{itemize}
\item Mesurer la même répartition des infaisables pour la machine de factorisation et le modèle bayésien : la part d'échecs d'XFOIL de ce dernier (41~\%) suggère que le modèle d'objectif influe aussi sur l'endroit où l'on propose.
\end{itemize}
\end{pistes}

\chapter{Classer les contraintes avant de les traiter}\label{ch:taxonomie}
\iterbanner{V3}

\section{Quatre questions}
En optimisation par simulation, la référence est la taxonomie de Le Digabel et Wild (\emph{A taxonomy of constraints in black-box simulation-based optimization}, Optimization and Engineering, 2023). Elle range chaque contrainte en répondant à quatre questions :
\begin{description}
\item[Connue ou cachée ?] Une contrainte \emph{connue} est écrite dans la formulation ; une contrainte \emph{cachée} n'apparaît que parce que le simulateur échoue (maillage impossible, solveur qui diverge).
\item[A priori ou issue de la simulation ?] Une contrainte \emph{a priori} se vérifie sans simuler (une borne, une géométrie) ; les autres demandent le calcul.
\item[Quantifiable ou non ?] Sait-on \emph{de combien} elle est violée, ou seulement \emph{si} elle l'est ?
\item[Relâchable ou non ?] Le simulateur donne-t-il encore des sorties utiles quand elle est violée ? Une contrainte non relâchable (une géométrie qui se croise) rend la simulation impossible ou dénuée de sens.
\end{description}
Une contrainte cachée est la pire : non quantifiable, non relâchable, issue de la simulation. Les auteurs la décrivent comme « du logiciel qui plante ».

\section{Nos quatre contraintes}
\begin{center}
\renewcommand{\arraystretch}{1.25}
\small
\begin{tabularx}{\linewidth}{>{\raggedright\arraybackslash}p{3.1cm} >{\raggedright\arraybackslash}p{3.8cm} >{\raggedright\arraybackslash}X}
\toprule
contrainte & classe & traitement adapté \\
\midrule
faces sans croisement & connue, a priori, quantifiable, \textbf{non relâchable} & ne jamais simuler un profil qui la viole : filtre avant simulation (« barrière extrême ») \\
épaisseur $\ge8~\%$ & connue, a priori, quantifiable, relâchable & filtre, ou formulation qui l'impose \\
$C_{L,\max}\ge1{,}2$ & connue, \textbf{issue de la simulation}, quantifiable, relâchable & un modèle de la contrainte, appris à part, avec son incertitude \\
convergence d'XFOIL & \textbf{cachée}, non quantifiable, non relâchable & une probabilité de succès apprise (classifieur) \\
\bottomrule
\end{tabularx}
\end{center}
La boucle de la V2 les traitait toutes par une pénalité quadratique, sauf la convergence, qu'elle ignorait jusqu'à la simulation. Le lot L6 a ajouté le filtre géométrique, qui est la bonne réponse pour les deux premières.

\section{Ce que fait la pratique industrielle}
Les chaînes d'optimisation de conception (aérodynamique, structures, moteurs) suivent la taxonomie, plus ou moins explicitement :
\begin{itemize}
\item \textbf{Contraintes a priori} : vérifiées avant de lancer un calcul, ou rendues impossibles à violer par la paramétrisation (on paramètre une loi d'épaisseur positive plutôt que deux faces indépendantes).
\item \textbf{Contraintes issues de la simulation} : un modèle de substitution \emph{par contrainte}, à côté de ceux des objectifs. En optimisation bayésienne contrainte (méthode SEGO, outils SMT et SEGOMOE de l'ONERA, utilisés dans le lot L1), le critère d'acquisition est multiplié par la \emph{probabilité d'être faisable} prédite. Une contrainte n'est jamais mélangée à l'objectif.
\item \textbf{Contraintes cachées} : un classifieur sépare les points où le code a abouti des autres (Gramacy et Lee, 2010 ; Gelbart et al., 2014), et sa probabilité de succès pondère le critère, comme celle d'une contrainte ordinaire.
\item \textbf{Solveurs} : un solveur industriel en variables binaires ou mixtes (CPLEX, Gurobi, SCIP, ou le solveur hybride contraint de D-Wave) reçoit l'objectif \emph{et} la liste des contraintes, et c'est lui qui choisit comment les traiter (chapitre~\ref{ch:cqm}).
\end{itemize}
Le QUBO est l'exception : il n'admet pas de contrainte, donc tout doit devenir un terme d'énergie. La question du cours est de savoir \emph{quels} termes, et ce qui peut rester en dehors du QUBO.

\begin{projet}
Une machine quantique, recuit ou circuit, rend une \emph{liste d'échantillons}, pas une seule réponse. On peut toujours les vérifier en classique avant de simuler : c'est déjà ce que fait le filtre géométrique. Cela ouvre une troisième voie entre « tout dans le QUBO » et « rien dans le QUBO » : mettre dans le QUBO ce qui \emph{oriente} la recherche, et laisser à la vérification ce qui \emph{garantit} la faisabilité.
\end{projet}

\begin{piege}[ : relâchable ne veut pas dire « peu importante »]
La contrainte de manœuvre est relâchable (XFOIL calcule quand même un profil qui ne tient pas $C_{L,\max}=1{,}2$), mais elle est dure pour le résultat : un profil qui la viole est inutilisable. « Relâchable » dit seulement qu'on peut apprendre quelque chose d'un point infaisable, ce qui est précieux pour un modèle de contrainte.
\end{piege}

\begin{resultat}
\begin{itemize}
\item On classe une contrainte par quatre questions : connue ou cachée, a priori ou simulée, quantifiable ou non, relâchable ou non.
\item Contrainte a priori : on filtre. Contrainte simulée quantifiable : on la modélise à part. Contrainte cachée : on apprend une probabilité de succès.
\item Les solveurs industriels reçoivent les contraintes séparément de l'objectif. Le QUBO ne le permet pas, mais ses échantillons peuvent être vérifiés en classique.
\end{itemize}
\end{resultat}

\exercices
\begin{exercice}\label{ex:taxo-classer}
Classer selon les quatre questions : (a) l'envergure d'une aile, bornée à 2~m ; (b) un mailleur qui échoue sur certaines géométries ; (c) une contrainte de flutter calculée par un code aéroélastique, qui donne une marge de vitesse ; (d) un calcul RANS qui converge mais vers une solution instationnaire, détectée par l'oscillation des résidus. Proposer un traitement pour chacune.
\end{exercice}
\begin{exercice}\label{ex:taxo-filtre}
Pourquoi vaut-il mieux filtrer une contrainte a priori non relâchable que la pénaliser, même lourdement ? Donner un argument sur le budget de simulations et un argument sur le modèle de substitution.
\end{exercice}

\begin{pistes}
\begin{itemize}
\item Enregistrer la \emph{cause} des échecs d'XFOIL (non-convergence, décollement massif, traînée aberrante) : un indicateur d'échec plus riche transforme une contrainte cachée en contrainte connue et non quantifiable, mieux modélisable.
\end{itemize}
\end{pistes}

\chapter{Ce que valent les pénalités}\label{ch:penalites}
\iterbanner{V3}

\section{La pénalité avec bits d'écart}
Pour une contrainte $h(\mathbf b)\ge0$ affine en bits, la formulation classique introduit un écart $s\ge0$ et pénalise l'égalité $h=s$ :
\[ P\,\big(h(\mathbf b)-s\big)^2,\qquad s=\frac{s_{\max}}{2^{n_s}-1}\sum_{j=0}^{n_s-1}2^j\,a_j,\quad a_j\in\{0,1\}. \]
Au minimum sur les bits d'écart $a_j$, ce terme vaut $P\,h^2$ si $h<0$ (on paie le carré de la violation) et $P\,(h-s^\star)^2$ si $h\ge0$, où $s^\star$ est le niveau d'écart le plus proche de $h$. Dans notre boucle : $P=75$, 3 bits d'écart, $s_{\max}$ égal à la plus grande marge possible.

\paragraph{Quatre défauts, tous observés.}
\begin{enumerate}
\item \textbf{Une petite violation coûte presque rien.} Une violation de 0,05 sur $C_{L,\max}$ coûte $75\times0{,}05^2\approx0{,}19$, alors que l'objectif normalisé s'étend sur environ 1. Si violer la contrainte fait gagner plus de 0,19 sur l'objectif, l'optimum du QUBO la viole. C'est ce qu'on mesure au chapitre~\ref{ch:hors-ligne} : le candidat retenu a un $\hat C_{L,\max}$ médian de 1,20, \emph{pile} sur la borne, et viole la contrainte du modèle une fois sur deux. Une pénalité quadratique n'est exacte que pour $P\to\infty$ ; une pénalité exacte serait en $|h|$, qui n'est pas quadratique.
\item \textbf{Les dents de scie.} Pour $h\ge0$, l'écart ne prend que 8 valeurs : le terme vaut $P(h-s^\star)^2$, nul seulement aux niveaux de l'écart, jusqu'à $P\Delta^2/4$ entre eux ($\Delta=s_{\max}/7$). Des profils également faisables reçoivent des pénalités différentes, sans raison physique (figure~\ref{fig:penalites}).
\item \textbf{Trois qubits de plus}, couplés à tous les bits de conception : 33 qubits au lieu de 30 à l'échelle, 93 couplages supplémentaires.
\item \textbf{Le spectre est écrasé.} Pour que la pénalité tienne, $P$ doit être grand ; l'objectif n'occupe alors que 1,7~\% (V2) ou 0,8~\% (CST) de l'étendue des énergies. Pour un recuit simulé réglé sur l'objectif, c'est gérable ; pour QAOA ou un recuit quantique, dont la précision est relative à l'étendue totale, c'est rédhibitoire.
\end{enumerate}
Il y a un cinquième défaut, de principe : le carré d'une contrainte quadratique est de degré 4. On ne peut donc pénaliser ainsi qu'une contrainte \emph{affine}, ce qui exclut un modèle quadratique de $C_{L,\max}$ et le classifieur du chapitre~\ref{ch:classifieur}.

\begin{figure}[htbp]
\centering
% figure générée par tools/make_figures.py
\begin{tikzpicture}
\begin{axis}[width=0.92\linewidth,height=6.2cm,xmin=-0.3,xmax=0.5,ymin=-0.6,ymax=2.6,
  xlabel={marge $h=\hat C_{L,\max}-1{,}2$},ylabel={terme ajouté à l'énergie},
  grid=major,grid style={cmuted!20},legend cell align=left,legend style={font=\footnotesize,at={(0.5,-0.3)},anchor=north,legend columns=1},
  extra x ticks={0},extra x tick style={grid style={cink,thick}}]
\fill[corange!8] (axis cs:-0.3,-0.6) rectangle (axis cs:0,2.6);
\node[font=\footnotesize,cink] at (axis cs:-0.15,2.35) {violée};
\node[font=\footnotesize,cink] at (axis cs:0.25,2.35) {respectée};
\addplot[cblue,very thick] coordinates {(-0.3,6.75) (-0.298,6.6603) (-0.296,6.5712) (-0.294,6.4827) (-0.292,6.3948) (-0.29,6.3075) (-0.288,6.2208) (-0.286,6.1347) (-0.284,6.0492) (-0.282,5.9643) (-0.28,5.88) (-0.278,5.7963) (-0.276,5.7132) (-0.274,5.6307) (-0.272,5.5488) (-0.27,5.4675) (-0.268,5.3868) (-0.266,5.3067) (-0.264,5.2272) (-0.262,5.1483) (-0.26,5.07) (-0.258,4.9923) (-0.256,4.9152) (-0.254,4.8387) (-0.252,4.7628) (-0.25,4.6875) (-0.248,4.6128) (-0.246,4.5387) (-0.244,4.4652) (-0.242,4.3923) (-0.24,4.32) (-0.238,4.2483) (-0.236,4.1772) (-0.234,4.1067) (-0.232,4.0368) (-0.23,3.9675) (-0.228,3.8988) (-0.226,3.8307) (-0.224,3.7632) (-0.222,3.6963) (-0.22,3.63) (-0.218,3.5643) (-0.216,3.4992) (-0.214,3.4347) (-0.212,3.3708) (-0.21,3.3075) (-0.208,3.2448) (-0.206,3.1827) (-0.204,3.1212) (-0.202,3.0603) (-0.2,3) (-0.198,2.9403) (-0.196,2.8812) (-0.194,2.8227) (-0.192,2.7648) (-0.19,2.7075) (-0.188,2.6508) (-0.186,2.5947) (-0.184,2.5392) (-0.182,2.4843) (-0.18,2.43) (-0.178,2.3763) (-0.176,2.3232) (-0.174,2.2707) (-0.172,2.2188) (-0.17,2.1675) (-0.168,2.1168) (-0.166,2.0667) (-0.164,2.0172) (-0.162,1.9683) (-0.16,1.92) (-0.158,1.8723) (-0.156,1.8252) (-0.154,1.7787) (-0.152,1.7328) (-0.15,1.6875) (-0.148,1.6428) (-0.146,1.5987) (-0.144,1.5552) (-0.142,1.5123) (-0.14,1.47) (-0.138,1.4283) (-0.136,1.3872) (-0.134,1.3467) (-0.132,1.3068) (-0.13,1.2675) (-0.128,1.2288) (-0.126,1.1907) (-0.124,1.1532) (-0.122,1.1163) (-0.12,1.08) (-0.118,1.0443) (-0.116,1.0092) (-0.114,0.9747) (-0.112,0.9408) (-0.11,0.9075) (-0.108,0.8748) (-0.106,0.8427) (-0.104,0.8112) (-0.102,0.7803) (-0.1,0.75) (-0.098,0.7203) (-0.096,0.6912) (-0.094,0.6627) (-0.092,0.6348) (-0.09,0.6075) (-0.088,0.5808) (-0.086,0.5547) (-0.084,0.5292) (-0.082,0.5043) (-0.08,0.48) (-0.078,0.4563) (-0.076,0.4332) (-0.074,0.4107) (-0.072,0.3888) (-0.07,0.3675) (-0.068,0.3468) (-0.066,0.3267) (-0.064,0.3072) (-0.062,0.2883) (-0.06,0.27) (-0.058,0.2523) (-0.056,0.2352) (-0.054,0.2187) (-0.052,0.2028) (-0.05,0.1875) (-0.048,0.1728) (-0.046,0.1587) (-0.044,0.1452) (-0.042,0.1323) (-0.04,0.12) (-0.038,0.1083) (-0.036,0.0972) (-0.034,0.0867) (-0.032,0.0768) (-0.03,0.0675) (-0.028,0.0588) (-0.026,0.0507) (-0.024,0.0432) (-0.022,0.0363) (-0.02,0.03) (-0.018,0.0243) (-0.016,0.0192) (-0.014,0.0147) (-0.012,0.0108) (-0.01,0.0075) (-0.008,0.0048) (-0.006,0.0027) (-0.004,0.0012) (-0.002,0.0003) (0,0) (0.002,0.0003) (0.004,0.0012) (0.006,0.0027) (0.008,0.0048) (0.01,0.0075) (0.012,0.0108) (0.014,0.0147) (0.016,0.0192) (0.018,0.0243) (0.02,0.03) (0.022,0.0363) (0.024,0.0432) (0.026,0.0507) (0.028,0.0588) (0.03,0.0675) (0.032,0.0768) (0.034,0.0867) (0.036,0.094139) (0.038,0.08381) (0.04,0.074082) (0.042,0.064953) (0.044,0.056424) (0.046,0.048496) (0.048,0.041167) (0.05,0.034439) (0.052,0.02831) (0.054,0.022782) (0.056,0.017853) (0.058,0.013524) (0.06,0.0097959) (0.062,0.0066673) (0.064,0.0041388) (0.066,0.0022102) (0.068,0.00088163) (0.07,0.00015306) (0.072,2.449e-05) (0.074,0.00049592) (0.076,0.0015673) (0.078,0.0032388) (0.08,0.0055102) (0.082,0.0083816) (0.084,0.011853) (0.086,0.015924) (0.088,0.020596) (0.09,0.025867) (0.092,0.031739) (0.094,0.03821) (0.096,0.045282) (0.098,0.052953) (0.1,0.061224) (0.102,0.070096) (0.104,0.079567) (0.106,0.089639) (0.108,0.091127) (0.11,0.080969) (0.112,0.071412) (0.114,0.062455) (0.116,0.054098) (0.118,0.046341) (0.12,0.039184) (0.122,0.032627) (0.124,0.026669) (0.126,0.021312) (0.128,0.016555) (0.13,0.012398) (0.132,0.0088408) (0.134,0.0058837) (0.136,0.0035265) (0.138,0.0017694) (0.14,0.00061224) (0.142,5.5102e-05) (0.144,9.7959e-05) (0.146,0.00074082) (0.148,0.0019837) (0.15,0.0038265) (0.152,0.0062694) (0.154,0.0093122) (0.156,0.012955) (0.158,0.017198) (0.16,0.022041) (0.162,0.027484) (0.164,0.033527) (0.166,0.040169) (0.168,0.047412) (0.17,0.055255) (0.172,0.063698) (0.174,0.072741) (0.176,0.082384) (0.178,0.092627) (0.18,0.088163) (0.182,0.078178) (0.184,0.068792) (0.186,0.060006) (0.188,0.05182) (0.19,0.044235) (0.192,0.037249) (0.194,0.030863) (0.196,0.025078) (0.198,0.019892) (0.2,0.015306) (0.202,0.01132) (0.204,0.0079347) (0.206,0.005149) (0.208,0.0029633) (0.21,0.0013776) (0.212,0.00039184) (0.214,6.1224e-06) (0.216,0.00022041) (0.218,0.0010347) (0.22,0.002449) (0.222,0.0044633) (0.224,0.0070776) (0.226,0.010292) (0.228,0.014106) (0.23,0.01852) (0.232,0.023535) (0.234,0.029149) (0.236,0.035363) (0.238,0.042178) (0.24,0.049592) (0.242,0.057606) (0.244,0.06622) (0.246,0.075435) (0.248,0.085249) (0.25,0.095663) (0.252,0.085249) (0.254,0.075435) (0.256,0.06622) (0.258,0.057606) (0.26,0.049592) (0.262,0.042178) (0.264,0.035363) (0.266,0.029149) (0.268,0.023535) (0.27,0.01852) (0.272,0.014106) (0.274,0.010292) (0.276,0.0070776) (0.278,0.0044633) (0.28,0.002449) (0.282,0.0010347) (0.284,0.00022041) (0.286,6.1224e-06) (0.288,0.00039184) (0.29,0.0013776) (0.292,0.0029633) (0.294,0.005149) (0.296,0.0079347) (0.298,0.01132) (0.3,0.015306) (0.302,0.019892) (0.304,0.025078) (0.306,0.030863) (0.308,0.037249) (0.31,0.044235) (0.312,0.05182) (0.314,0.060006) (0.316,0.068792) (0.318,0.078178) (0.32,0.088163) (0.322,0.092627) (0.324,0.082384) (0.326,0.072741) (0.328,0.063698) (0.33,0.055255) (0.332,0.047412) (0.334,0.040169) (0.336,0.033527) (0.338,0.027484) (0.34,0.022041) (0.342,0.017198) (0.344,0.012955) (0.346,0.0093122) (0.348,0.0062694) (0.35,0.0038265) (0.352,0.0019837) (0.354,0.00074082) (0.356,9.7959e-05) (0.358,5.5102e-05) (0.36,0.00061224) (0.362,0.0017694) (0.364,0.0035265) (0.366,0.0058837) (0.368,0.0088408) (0.37,0.012398) (0.372,0.016555) (0.374,0.021312) (0.376,0.026669) (0.378,0.032627) (0.38,0.039184) (0.382,0.046341) (0.384,0.054098) (0.386,0.062455) (0.388,0.071412) (0.39,0.080969) (0.392,0.091127) (0.394,0.089639) (0.396,0.079567) (0.398,0.070096) (0.4,0.061224) (0.402,0.052953) (0.404,0.045282) (0.406,0.03821) (0.408,0.031739) (0.41,0.025867) (0.412,0.020596) (0.414,0.015924) (0.416,0.011853) (0.418,0.0083816) (0.42,0.0055102) (0.422,0.0032388) (0.424,0.0015673) (0.426,0.00049592) (0.428,2.449e-05) (0.43,0.00015306) (0.432,0.00088163) (0.434,0.0022102) (0.436,0.0041388) (0.438,0.0066673) (0.44,0.0097959) (0.442,0.013524) (0.444,0.017853) (0.446,0.022782) (0.448,0.02831) (0.45,0.034439) (0.452,0.041167) (0.454,0.048496) (0.456,0.056424) (0.458,0.064953) (0.46,0.074082) (0.462,0.08381) (0.464,0.094139) (0.466,0.0867) (0.468,0.0768) (0.47,0.0675) (0.472,0.0588) (0.474,0.0507) (0.476,0.0432) (0.478,0.0363) (0.48,0.03) (0.482,0.0243) (0.484,0.0192) (0.486,0.0147) (0.488,0.0108) (0.49,0.0075) (0.492,0.0048) (0.494,0.0027) (0.496,0.0012) (0.498,0.0003) (0.5,0)};
\addlegendentry{bits d'écart (V2) : $P\min_s(h-s)^2$, 3 bits}
\addplot[corange,very thick] coordinates {(-0.3,9) (-0.298,8.8953) (-0.296,8.7912) (-0.294,8.6877) (-0.292,8.5848) (-0.29,8.4825) (-0.288,8.3808) (-0.286,8.2797) (-0.284,8.1792) (-0.282,8.0793) (-0.28,7.98) (-0.278,7.8813) (-0.276,7.7832) (-0.274,7.6857) (-0.272,7.5888) (-0.27,7.4925) (-0.268,7.3968) (-0.266,7.3017) (-0.264,7.2072) (-0.262,7.1133) (-0.26,7.02) (-0.258,6.9273) (-0.256,6.8352) (-0.254,6.7437) (-0.252,6.6528) (-0.25,6.5625) (-0.248,6.4728) (-0.246,6.3837) (-0.244,6.2952) (-0.242,6.2073) (-0.24,6.12) (-0.238,6.0333) (-0.236,5.9472) (-0.234,5.8617) (-0.232,5.7768) (-0.23,5.6925) (-0.228,5.6088) (-0.226,5.5257) (-0.224,5.4432) (-0.222,5.3613) (-0.22,5.28) (-0.218,5.1993) (-0.216,5.1192) (-0.214,5.0397) (-0.212,4.9608) (-0.21,4.8825) (-0.208,4.8048) (-0.206,4.7277) (-0.204,4.6512) (-0.202,4.5753) (-0.2,4.5) (-0.198,4.4253) (-0.196,4.3512) (-0.194,4.2777) (-0.192,4.2048) (-0.19,4.1325) (-0.188,4.0608) (-0.186,3.9897) (-0.184,3.9192) (-0.182,3.8493) (-0.18,3.78) (-0.178,3.7113) (-0.176,3.6432) (-0.174,3.5757) (-0.172,3.5088) (-0.17,3.4425) (-0.168,3.3768) (-0.166,3.3117) (-0.164,3.2472) (-0.162,3.1833) (-0.16,3.12) (-0.158,3.0573) (-0.156,2.9952) (-0.154,2.9337) (-0.152,2.8728) (-0.15,2.8125) (-0.148,2.7528) (-0.146,2.6937) (-0.144,2.6352) (-0.142,2.5773) (-0.14,2.52) (-0.138,2.4633) (-0.136,2.4072) (-0.134,2.3517) (-0.132,2.2968) (-0.13,2.2425) (-0.128,2.1888) (-0.126,2.1357) (-0.124,2.0832) (-0.122,2.0313) (-0.12,1.98) (-0.118,1.9293) (-0.116,1.8792) (-0.114,1.8297) (-0.112,1.7808) (-0.11,1.7325) (-0.108,1.6848) (-0.106,1.6377) (-0.104,1.5912) (-0.102,1.5453) (-0.1,1.5) (-0.098,1.4553) (-0.096,1.4112) (-0.094,1.3677) (-0.092,1.3248) (-0.09,1.2825) (-0.088,1.2408) (-0.086,1.1997) (-0.084,1.1592) (-0.082,1.1193) (-0.08,1.08) (-0.078,1.0413) (-0.076,1.0032) (-0.074,0.9657) (-0.072,0.9288) (-0.07,0.8925) (-0.068,0.8568) (-0.066,0.8217) (-0.064,0.7872) (-0.062,0.7533) (-0.06,0.72) (-0.058,0.6873) (-0.056,0.6552) (-0.054,0.6237) (-0.052,0.5928) (-0.05,0.5625) (-0.048,0.5328) (-0.046,0.5037) (-0.044,0.4752) (-0.042,0.4473) (-0.04,0.42) (-0.038,0.3933) (-0.036,0.3672) (-0.034,0.3417) (-0.032,0.3168) (-0.03,0.2925) (-0.028,0.2688) (-0.026,0.2457) (-0.024,0.2232) (-0.022,0.2013) (-0.02,0.18) (-0.018,0.1593) (-0.016,0.1392) (-0.014,0.1197) (-0.012,0.1008) (-0.01,0.0825) (-0.008,0.0648) (-0.006,0.0477) (-0.004,0.0312) (-0.002,0.0153) (0,0) (0.002,-0.0147) (0.004,-0.0288) (0.006,-0.0423) (0.008,-0.0552) (0.01,-0.0675) (0.012,-0.0792) (0.014,-0.0903) (0.016,-0.1008) (0.018,-0.1107) (0.02,-0.12) (0.022,-0.1287) (0.024,-0.1368) (0.026,-0.1443) (0.028,-0.1512) (0.03,-0.1575) (0.032,-0.1632) (0.034,-0.1683) (0.036,-0.1728) (0.038,-0.1767) (0.04,-0.18) (0.042,-0.1827) (0.044,-0.1848) (0.046,-0.1863) (0.048,-0.1872) (0.05,-0.1875) (0.052,-0.1872) (0.054,-0.1863) (0.056,-0.1848) (0.058,-0.1827) (0.06,-0.18) (0.062,-0.1767) (0.064,-0.1728) (0.066,-0.1683) (0.068,-0.1632) (0.07,-0.1575) (0.072,-0.1512) (0.074,-0.1443) (0.076,-0.1368) (0.078,-0.1287) (0.08,-0.12) (0.082,-0.1107) (0.084,-0.1008) (0.086,-0.0903) (0.088,-0.0792) (0.09,-0.0675) (0.092,-0.0552) (0.094,-0.0423) (0.096,-0.0288) (0.098,-0.0147) (0.1,2.6021e-16) (0.102,0.0153) (0.104,0.0312) (0.106,0.0477) (0.108,0.0648) (0.11,0.0825) (0.112,0.1008) (0.114,0.1197) (0.116,0.1392) (0.118,0.1593) (0.12,0.18) (0.122,0.2013) (0.124,0.2232) (0.126,0.2457) (0.128,0.2688) (0.13,0.2925) (0.132,0.3168) (0.134,0.3417) (0.136,0.3672) (0.138,0.3933) (0.14,0.42) (0.142,0.4473) (0.144,0.4752) (0.146,0.5037) (0.148,0.5328) (0.15,0.5625) (0.152,0.5928) (0.154,0.6237) (0.156,0.6552) (0.158,0.6873) (0.16,0.72) (0.162,0.7533) (0.164,0.7872) (0.166,0.8217) (0.168,0.8568) (0.17,0.8925) (0.172,0.9288) (0.174,0.9657) (0.176,1.0032) (0.178,1.0413) (0.18,1.08) (0.182,1.1193) (0.184,1.1592) (0.186,1.1997) (0.188,1.2408) (0.19,1.2825) (0.192,1.3248) (0.194,1.3677) (0.196,1.4112) (0.198,1.4553) (0.2,1.5) (0.202,1.5453) (0.204,1.5912) (0.206,1.6377) (0.208,1.6848) (0.21,1.7325) (0.212,1.7808) (0.214,1.8297) (0.216,1.8792) (0.218,1.9293) (0.22,1.98) (0.222,2.0313) (0.224,2.0832) (0.226,2.1357) (0.228,2.1888) (0.23,2.2425) (0.232,2.2968) (0.234,2.3517) (0.236,2.4072) (0.238,2.4633) (0.24,2.52) (0.242,2.5773) (0.244,2.6352) (0.246,2.6937) (0.248,2.7528) (0.25,2.8125) (0.252,2.8728) (0.254,2.9337) (0.256,2.9952) (0.258,3.0573) (0.26,3.12) (0.262,3.1833) (0.264,3.2472) (0.266,3.3117) (0.268,3.3768) (0.27,3.4425) (0.272,3.5088) (0.274,3.5757) (0.276,3.6432) (0.278,3.7113) (0.28,3.78) (0.282,3.8493) (0.284,3.9192) (0.286,3.9897) (0.288,4.0608) (0.29,4.1325) (0.292,4.2048) (0.294,4.2777) (0.296,4.3512) (0.298,4.4253) (0.3,4.5) (0.302,4.5753) (0.304,4.6512) (0.306,4.7277) (0.308,4.8048) (0.31,4.8825) (0.312,4.9608) (0.314,5.0397) (0.316,5.1192) (0.318,5.1993) (0.32,5.28) (0.322,5.3613) (0.324,5.4432) (0.326,5.5257) (0.328,5.6088) (0.33,5.6925) (0.332,5.7768) (0.334,5.8617) (0.336,5.9472) (0.338,6.0333) (0.34,6.12) (0.342,6.2073) (0.344,6.2952) (0.346,6.3837) (0.348,6.4728) (0.35,6.5625) (0.352,6.6528) (0.354,6.7437) (0.356,6.8352) (0.358,6.9273) (0.36,7.02) (0.362,7.1133) (0.364,7.2072) (0.366,7.3017) (0.368,7.3968) (0.37,7.4925) (0.372,7.5888) (0.374,7.6857) (0.376,7.7832) (0.378,7.8813) (0.38,7.98) (0.382,8.0793) (0.384,8.1792) (0.386,8.2797) (0.388,8.3808) (0.39,8.4825) (0.392,8.5848) (0.394,8.6877) (0.396,8.7912) (0.398,8.8953) (0.4,9) (0.402,9.1053) (0.404,9.2112) (0.406,9.3177) (0.408,9.4248) (0.41,9.5325) (0.412,9.6408) (0.414,9.7497) (0.416,9.8592) (0.418,9.9693) (0.42,10.08) (0.422,10.191) (0.424,10.303) (0.426,10.416) (0.428,10.529) (0.43,10.642) (0.432,10.757) (0.434,10.872) (0.436,10.987) (0.438,11.103) (0.44,11.22) (0.442,11.337) (0.444,11.455) (0.446,11.574) (0.448,11.693) (0.45,11.812) (0.452,11.933) (0.454,12.054) (0.456,12.175) (0.458,12.297) (0.46,12.42) (0.462,12.543) (0.464,12.667) (0.466,12.792) (0.468,12.917) (0.47,13.043) (0.472,13.169) (0.474,13.296) (0.476,13.423) (0.478,13.551) (0.48,13.68) (0.482,13.809) (0.484,13.939) (0.486,14.07) (0.488,14.201) (0.49,14.333) (0.492,14.465) (0.494,14.598) (0.496,14.731) (0.498,14.865) (0.5,15)};
\addlegendentry{déséquilibrée : $P(h^2-2h^\star h)$, $h^\star=0{,}05$}
\addplot[cgreen,very thick] coordinates {(-0.3,0.24) (-0.298,0.2384) (-0.296,0.2368) (-0.294,0.2352) (-0.292,0.2336) (-0.29,0.232) (-0.288,0.2304) (-0.286,0.2288) (-0.284,0.2272) (-0.282,0.2256) (-0.28,0.224) (-0.278,0.2224) (-0.276,0.2208) (-0.274,0.2192) (-0.272,0.2176) (-0.27,0.216) (-0.268,0.2144) (-0.266,0.2128) (-0.264,0.2112) (-0.262,0.2096) (-0.26,0.208) (-0.258,0.2064) (-0.256,0.2048) (-0.254,0.2032) (-0.252,0.2016) (-0.25,0.2) (-0.248,0.1984) (-0.246,0.1968) (-0.244,0.1952) (-0.242,0.1936) (-0.24,0.192) (-0.238,0.1904) (-0.236,0.1888) (-0.234,0.1872) (-0.232,0.1856) (-0.23,0.184) (-0.228,0.1824) (-0.226,0.1808) (-0.224,0.1792) (-0.222,0.1776) (-0.22,0.176) (-0.218,0.1744) (-0.216,0.1728) (-0.214,0.1712) (-0.212,0.1696) (-0.21,0.168) (-0.208,0.1664) (-0.206,0.1648) (-0.204,0.1632) (-0.202,0.1616) (-0.2,0.16) (-0.198,0.1584) (-0.196,0.1568) (-0.194,0.1552) (-0.192,0.1536) (-0.19,0.152) (-0.188,0.1504) (-0.186,0.1488) (-0.184,0.1472) (-0.182,0.1456) (-0.18,0.144) (-0.178,0.1424) (-0.176,0.1408) (-0.174,0.1392) (-0.172,0.1376) (-0.17,0.136) (-0.168,0.1344) (-0.166,0.1328) (-0.164,0.1312) (-0.162,0.1296) (-0.16,0.128) (-0.158,0.1264) (-0.156,0.1248) (-0.154,0.1232) (-0.152,0.1216) (-0.15,0.12) (-0.148,0.1184) (-0.146,0.1168) (-0.144,0.1152) (-0.142,0.1136) (-0.14,0.112) (-0.138,0.1104) (-0.136,0.1088) (-0.134,0.1072) (-0.132,0.1056) (-0.13,0.104) (-0.128,0.1024) (-0.126,0.1008) (-0.124,0.0992) (-0.122,0.0976) (-0.12,0.096) (-0.118,0.0944) (-0.116,0.0928) (-0.114,0.0912) (-0.112,0.0896) (-0.11,0.088) (-0.108,0.0864) (-0.106,0.0848) (-0.104,0.0832) (-0.102,0.0816) (-0.1,0.08) (-0.098,0.0784) (-0.096,0.0768) (-0.094,0.0752) (-0.092,0.0736) (-0.09,0.072) (-0.088,0.0704) (-0.086,0.0688) (-0.084,0.0672) (-0.082,0.0656) (-0.08,0.064) (-0.078,0.0624) (-0.076,0.0608) (-0.074,0.0592) (-0.072,0.0576) (-0.07,0.056) (-0.068,0.0544) (-0.066,0.0528) (-0.064,0.0512) (-0.062,0.0496) (-0.06,0.048) (-0.058,0.0464) (-0.056,0.0448) (-0.054,0.0432) (-0.052,0.0416) (-0.05,0.04) (-0.048,0.0384) (-0.046,0.0368) (-0.044,0.0352) (-0.042,0.0336) (-0.04,0.032) (-0.038,0.0304) (-0.036,0.0288) (-0.034,0.0272) (-0.032,0.0256) (-0.03,0.024) (-0.028,0.0224) (-0.026,0.0208) (-0.024,0.0192) (-0.022,0.0176) (-0.02,0.016) (-0.018,0.0144) (-0.016,0.0128) (-0.014,0.0112) (-0.012,0.0096) (-0.01,0.008) (-0.008,0.0064) (-0.006,0.0048) (-0.004,0.0032) (-0.002,0.0016) (0,-0) (0.002,-0.0016) (0.004,-0.0032) (0.006,-0.0048) (0.008,-0.0064) (0.01,-0.008) (0.012,-0.0096) (0.014,-0.0112) (0.016,-0.0128) (0.018,-0.0144) (0.02,-0.016) (0.022,-0.0176) (0.024,-0.0192) (0.026,-0.0208) (0.028,-0.0224) (0.03,-0.024) (0.032,-0.0256) (0.034,-0.0272) (0.036,-0.0288) (0.038,-0.0304) (0.04,-0.032) (0.042,-0.0336) (0.044,-0.0352) (0.046,-0.0368) (0.048,-0.0384) (0.05,-0.04) (0.052,-0.0416) (0.054,-0.0432) (0.056,-0.0448) (0.058,-0.0464) (0.06,-0.048) (0.062,-0.0496) (0.064,-0.0512) (0.066,-0.0528) (0.068,-0.0544) (0.07,-0.056) (0.072,-0.0576) (0.074,-0.0592) (0.076,-0.0608) (0.078,-0.0624) (0.08,-0.064) (0.082,-0.0656) (0.084,-0.0672) (0.086,-0.0688) (0.088,-0.0704) (0.09,-0.072) (0.092,-0.0736) (0.094,-0.0752) (0.096,-0.0768) (0.098,-0.0784) (0.1,-0.08) (0.102,-0.0816) (0.104,-0.0832) (0.106,-0.0848) (0.108,-0.0864) (0.11,-0.088) (0.112,-0.0896) (0.114,-0.0912) (0.116,-0.0928) (0.118,-0.0944) (0.12,-0.096) (0.122,-0.0976) (0.124,-0.0992) (0.126,-0.1008) (0.128,-0.1024) (0.13,-0.104) (0.132,-0.1056) (0.134,-0.1072) (0.136,-0.1088) (0.138,-0.1104) (0.14,-0.112) (0.142,-0.1136) (0.144,-0.1152) (0.146,-0.1168) (0.148,-0.1184) (0.15,-0.12) (0.152,-0.1216) (0.154,-0.1232) (0.156,-0.1248) (0.158,-0.1264) (0.16,-0.128) (0.162,-0.1296) (0.164,-0.1312) (0.166,-0.1328) (0.168,-0.1344) (0.17,-0.136) (0.172,-0.1376) (0.174,-0.1392) (0.176,-0.1408) (0.178,-0.1424) (0.18,-0.144) (0.182,-0.1456) (0.184,-0.1472) (0.186,-0.1488) (0.188,-0.1504) (0.19,-0.152) (0.192,-0.1536) (0.194,-0.1552) (0.196,-0.1568) (0.198,-0.1584) (0.2,-0.16) (0.202,-0.1616) (0.204,-0.1632) (0.206,-0.1648) (0.208,-0.1664) (0.21,-0.168) (0.212,-0.1696) (0.214,-0.1712) (0.216,-0.1728) (0.218,-0.1744) (0.22,-0.176) (0.222,-0.1776) (0.224,-0.1792) (0.226,-0.1808) (0.228,-0.1824) (0.23,-0.184) (0.232,-0.1856) (0.234,-0.1872) (0.236,-0.1888) (0.238,-0.1904) (0.24,-0.192) (0.242,-0.1936) (0.244,-0.1952) (0.246,-0.1968) (0.248,-0.1984) (0.25,-0.2) (0.252,-0.2016) (0.254,-0.2032) (0.256,-0.2048) (0.258,-0.2064) (0.26,-0.208) (0.262,-0.2096) (0.264,-0.2112) (0.266,-0.2128) (0.268,-0.2144) (0.27,-0.216) (0.272,-0.2176) (0.274,-0.2192) (0.276,-0.2208) (0.278,-0.2224) (0.28,-0.224) (0.282,-0.2256) (0.284,-0.2272) (0.286,-0.2288) (0.288,-0.2304) (0.29,-0.232) (0.292,-0.2336) (0.294,-0.2352) (0.296,-0.2368) (0.298,-0.2384) (0.3,-0.24) (0.302,-0.2416) (0.304,-0.2432) (0.306,-0.2448) (0.308,-0.2464) (0.31,-0.248) (0.312,-0.2496) (0.314,-0.2512) (0.316,-0.2528) (0.318,-0.2544) (0.32,-0.256) (0.322,-0.2576) (0.324,-0.2592) (0.326,-0.2608) (0.328,-0.2624) (0.33,-0.264) (0.332,-0.2656) (0.334,-0.2672) (0.336,-0.2688) (0.338,-0.2704) (0.34,-0.272) (0.342,-0.2736) (0.344,-0.2752) (0.346,-0.2768) (0.348,-0.2784) (0.35,-0.28) (0.352,-0.2816) (0.354,-0.2832) (0.356,-0.2848) (0.358,-0.2864) (0.36,-0.288) (0.362,-0.2896) (0.364,-0.2912) (0.366,-0.2928) (0.368,-0.2944) (0.37,-0.296) (0.372,-0.2976) (0.374,-0.2992) (0.376,-0.3008) (0.378,-0.3024) (0.38,-0.304) (0.382,-0.3056) (0.384,-0.3072) (0.386,-0.3088) (0.388,-0.3104) (0.39,-0.312) (0.392,-0.3136) (0.394,-0.3152) (0.396,-0.3168) (0.398,-0.3184) (0.4,-0.32) (0.402,-0.3216) (0.404,-0.3232) (0.406,-0.3248) (0.408,-0.3264) (0.41,-0.328) (0.412,-0.3296) (0.414,-0.3312) (0.416,-0.3328) (0.418,-0.3344) (0.42,-0.336) (0.422,-0.3376) (0.424,-0.3392) (0.426,-0.3408) (0.428,-0.3424) (0.43,-0.344) (0.432,-0.3456) (0.434,-0.3472) (0.436,-0.3488) (0.438,-0.3504) (0.44,-0.352) (0.442,-0.3536) (0.444,-0.3552) (0.446,-0.3568) (0.448,-0.3584) (0.45,-0.36) (0.452,-0.3616) (0.454,-0.3632) (0.456,-0.3648) (0.458,-0.3664) (0.46,-0.368) (0.462,-0.3696) (0.464,-0.3712) (0.466,-0.3728) (0.468,-0.3744) (0.47,-0.376) (0.472,-0.3776) (0.474,-0.3792) (0.476,-0.3808) (0.478,-0.3824) (0.48,-0.384) (0.482,-0.3856) (0.484,-0.3872) (0.486,-0.3888) (0.488,-0.3904) (0.49,-0.392) (0.492,-0.3936) (0.494,-0.3952) (0.496,-0.3968) (0.498,-0.3984) (0.5,-0.4)};
\addlegendentry{lagrangien : $-\lambda h/s$, $\lambda/s=0{,}8$}
\end{axis}
\end{tikzpicture}

\caption{Les trois termes ajoutés à l'énergie, en fonction de la marge $h$ ($P=75$, échelle de $C_{L,\max}$). Bits d'écart : violation payée au carré, dents de scie du côté faisable. Pénalité déséquilibrée : pas de qubit, mais minimum en $h^\star$ et croissance du côté faisable. Lagrangien : une droite, qui récompense la marge au lieu de la pénaliser.}\label{fig:penalites}
\end{figure}

\section{La pénalité déséquilibrée}
Montañez-Barrera, Willsch, Maldonado-Romo et Michielsen (\emph{Unbalanced penalization}, Quantum Science and Technology, 2024) suppriment les bits d'écart. L'idée part de la pénalité idéale $e^{-h}$, grande pour $h<0$ et presque nulle pour $h>0$, et en garde le développement de degré 2 : $1-h+h^2/2$. Le terme ajouté est $-\lambda_1h+\lambda_2h^2$, quadratique si $h$ est affine. Nous l'écrivons
\[ P\,\big(h^2-2h^\star h\big)=P\,\big[(h-h^\star)^2-h^{\star2}\big], \]
nulle en $h=0$, minimale en $h=h^\star>0$, positive pour $h<0$. Elle a été testée sur le voyageur de commerce, le sac à dos et le rangement en boîtes ; sur ce dernier, elle permet de traiter sur machine D-Wave des instances de 29 objets, contre 11 avec les bits d'écart.

\paragraph{Son défaut chez nous.} Elle pénalise aussi le côté faisable : avec $P=75$ et $h^\star=0{,}05$, un profil de marge $0{,}3$ paie $75\times(0{,}09-0{,}03)=4{,}5$, quatre fois l'étendue de l'objectif. Elle attire donc tous les candidats vers $h\approx h^\star$, ce qui déforme l'objectif (les « faux minimums » relevés dans la littérature). Sur des problèmes combinatoires où la contrainte est saturée à l'optimum (un sac à dos plein), c'est sans conséquence ; sur un front de Pareto dont une partie est loin de la contrainte, c'est un biais.

\section{Régler les poids : ce que dit la littérature}
Le réglage de $P$ est un sujet de recherche à lui seul. Les méthodes publiées se rangent en trois familles :
\begin{itemize}
\item \textbf{bornes exactes} : un $P$ assez grand pour qu'aucune violation ne soit rentable, calculé à partir des coefficients de l'objectif. Sûr, mais il écrase le spectre ;
\item \textbf{poids séquentiels} : partir d'un petit $P$ et l'augmenter tant que les échantillons sont infaisables (travaux de l'université d'Exeter, d'Ayodele) ;
\item \textbf{poids par contrainte} réglés sur le taux de faisabilité des échantillons, par exemple avec une loi sigmoïde.
\end{itemize}
Notre $P=75$, fixé à la main sur le problème V2, n'appartient à aucune. La deuxième famille est très proche du lagrangien du chapitre suivant, qui ajuste lui aussi un poids d'après les échantillons, mais sur un terme linéaire.

\begin{resultat}
\begin{itemize}
\item La pénalité quadratique avec bits d'écart rend les petites violations rentables, crée des dents de scie, ajoute 3 qubits et écrase le spectre. Elle ne s'applique qu'à une contrainte affine.
\item La pénalité déséquilibrée supprime les bits d'écart, mais attire tous les candidats vers une marge fixée : elle déforme l'objectif loin de la contrainte.
\item Toute pénalité au carré est limitée aux contraintes affines.
\end{itemize}
\end{resultat}

\exercices
\begin{exercice}\label{ex:pen-rentable}
Objectif normalisé : un profil faisable de marge $h=0$ vaut $-0{,}80$ ; son voisin vaut $-1{,}00$ avec $h=-0{,}04$. (a) Avec la pénalité à bits d'écart ($P=75$), lequel est l'optimum du QUBO ? (b) Quel $P$ faudrait-il pour que le profil faisable gagne ? (c) Que devient alors la part de l'objectif dans le spectre, si la marge prédite varie de $-0{,}5$ à $+0{,}5$ sur la grille ?
\end{exercice}
\begin{exercice}\label{ex:pen-deseq}
Pénalité déséquilibrée avec $P=75$, $h^\star=0{,}05$. (a) Calculer le terme en $h=-0{,}05$, $h=0$, $h=0{,}05$, $h=0{,}3$. (b) Pour quelles marges $h>0$ la pénalité est-elle positive ? (c) Proposer $P$ et $h^\star$ pour que le terme reste inférieur à 0,5 jusqu'à $h=0{,}3$, et dire ce que l'on perd.
\end{exercice}
\begin{exercice}\label{ex:pen-degre}
Le modèle de $C_{L,\max}$ devient quadratique en $z$. Montrer sur l'exemple $h=z_1z_2-0{,}2$ (deux variables à 1 bit, $z_i=2b_i-1$) que $h^2$ contient un monôme de degré supérieur à 2 en bits, sauf simplification due à $b_i^2=b_i$. Conclure pour des variables à 3 bits.
\end{exercice}

\begin{pistes}
\begin{itemize}
\item Régler $P$ par la méthode séquentielle (augmenter tant que la meilleure lecture viole la contrainte) au lieu de $P=75$ : c'est la comparaison la plus juste entre bits d'écart et lagrangien.
\end{itemize}
\end{pistes}

\chapter{Le lagrangien vérifié}\label{ch:lagrangien}
\iterbanner{V3}

\section{Un terme linéaire au lieu d'un carré}
Le problème que l'on veut résoudre à chaque itération, pour un poids $w$, est
\[ \min_{\mathbf b}\ f(\mathbf b)\quad\text{sous}\quad h_k(\mathbf b)\ge0,\ k=1,\dots,K, \]
où $f$ est l'objectif pondéré prédit et les $h_k$ les contraintes prédites, toutes écrites sous forme QUBO ($h_k(\mathbf b)=\mathbf b^\top H_k\mathbf b+h_{0,k}$). Le \textbf{lagrangien} remplace les contraintes par des termes linéaires :
\[ L(\mathbf b,\boldsymbol\lambda)=f(\mathbf b)-\sum_k\lambda_k\,\frac{h_k(\mathbf b)}{s_k},\qquad\lambda_k\ge0, \]
où $s_k$ est une échelle typique de $h_k$ (l'étendue de $C_{L,\max}$ sur les données, par exemple). Trois propriétés en font le bon outil pour un QUBO :
\begin{enumerate}
\item \textbf{le degré ne monte pas} : si $f$ et les $h_k$ sont quadratiques, $L$ l'est aussi. On peut donc garder une contrainte quadratique (un modèle quadratique de $C_{L,\max}$, le classifieur du chapitre~\ref{ch:classifieur}), ce qu'aucune pénalité au carré ne permet ;
\item \textbf{aucun qubit supplémentaire} : 30 qubits à l'échelle, ceux de la conception seule ;
\item \textbf{l'objectif garde le spectre} : avec des $\lambda_k$ petits, $L$ diffère peu de $f$, et la part de l'objectif dans l'étendue des énergies passe de 0,8~\% à environ 100~\% (chapitre~\ref{ch:hors-ligne}).
\end{enumerate}
Le prix : $-\lambda h$ \emph{récompense} la marge, sans limite. Le minimum de $L$ n'a aucune raison d'être sur la frontière ; avec des $\lambda$ trop grands, il part loin à l'intérieur du domaine faisable.

\section{Le saut de dualité, et pourquoi il faut vérifier}
En optimisation continue convexe, il existe des multiplicateurs $\boldsymbol\lambda^\star$ pour lesquels le minimum de $L(\cdot,\boldsymbol\lambda^\star)$ est exactement l'optimum contraint. En variables binaires, ce n'est plus vrai : il peut n'exister \emph{aucun} $\boldsymbol\lambda$ dont le minimum soit l'optimum contraint. C'est le \textbf{saut de dualité}. Le cas typique : l'optimum contraint est dominé, dans le plan (marge, objectif), par une combinaison d'un point très faisable et d'un point infaisable très bon (exercice~\ref{ex:lag-saut}).

La réponse est de ne pas demander au QUBO de rendre l'optimum, mais de bons \emph{échantillons} :
\begin{enumerate}
\item on résout le QUBO $L(\cdot,\boldsymbol\lambda)$ pour quelques valeurs de $\boldsymbol\lambda$ (6 résolutions) ;
\item après chaque résolution, on regarde la meilleure lecture : si elle viole la contrainte $k$, on double $\lambda_k$, sinon on le divise par deux (montée duale multiplicative), en partant de $\lambda_k=0{,}05$ ;
\item on réunit toutes les lectures de toutes les résolutions et l'on \textbf{vérifie} chacune en classique : on garde la meilleure, \emph{au sens de l'objectif seul}, parmi celles qui respectent toutes les contraintes prédites et le filtre géométrique.
\end{enumerate}
Les multiplicateurs servent à \emph{orienter} les lectures vers la frontière ; c'est la vérification qui garantit la faisabilité prédite et qui choisit selon l'objectif. Les $\lambda$ finaux d'une itération servent de départ à la suivante, pour le même poids.

\begin{codebox}[ : \code{contraintes/formulations.py}, \code{resout\_lagrange}]
\begin{lstlisting}
lam = np.full(K, lam_min)                  # 0,05 : presque l'objectif seul
for _ in range(n_tours):                   # 6 résolutions
    S, E = solve(*lagrangien(Qo, co, ctrs, lam))
    pool.append(S)
    m = filtre(S)                          # géométrie exacte
    best = S[m][np.argmin(E[m])]
    for k, ct in enumerate(ctrs):          # montée duale multiplicative
        lam[k] = 2 * lam[k] if ct.h(best) < 0 else lam[k] / 2
S = np.unique(np.vstack(pool), axis=0)     # vérification de toutes les lectures
ok = filtre(S) & all(ct.h(S) >= 0 for ct in ctrs if ct.verifiee)
ordre = faisables par objectif croissant, puis les autres par violation croissante
\end{lstlisting}
\end{codebox}

\begin{piege}[ : des multiplicateurs partis trop haut]
Une première version partait de $\lambda_k=1$. Avec sept contraintes (une de $C_{L,\max}$, six géométriques) dont les termes valent chacun 1 à 3, le lagrangien écrasait l'objectif : les candidats retenus avaient un $\hat C_{L,\max}$ médian de 1,53 et un objectif prédit moins bon de 0,95 (une étendue entière) que ceux des bits d'écart. Partir de $\lambda=0{,}05$ et ne monter qu'en cas de violation corrige ce défaut (chapitre~\ref{ch:hors-ligne}). Avec un terme linéaire, « trop contraint » n'est pas « plus sûr » : c'est un autre biais.
\end{piege}

\section{La géométrie est exactement affine en $z$}
Sans la géométrie dans le QUBO, le lagrangien retrouve le défaut de la première campagne CST : l'optimum du modèle est dans le coin des profils trop minces, dont les faces se croisent. Sur 40 QUBO de test, 13 ne donnaient \emph{aucune} lecture admissible. Or la géométrie s'écrit exactement. L'épaisseur à l'abscisse $x$ vaut
\[ e(x)=C(x)\sum_{i=0}^{4}B_i(x)\,\big(A_{{\rm ext},i}-A_{{\rm int},i}\big),\qquad A_j=\mathrm{LO}_j+\frac{\mathrm{HI}_j-\mathrm{LO}_j}{2}\,(z_j+1), \]
où $C(x)=\sqrt x(1-x)$ est la fonction de classe, $B_i$ les polynômes de Bernstein, et $A_j$ le coefficient $j$, affine en sa variable normalisée $z_j$. À $x$ fixé, $e(x)$ est donc \textbf{affine en $z$}, donc linéaire en bits (vérifié : écart de $10^{-16}$ avec la géométrie calculée).
\begin{itemize}
\item \textbf{Faces sans croisement} : $\sum_iB_i(x_j)(A_{{\rm ext},i}-A_{{\rm int},i})\ge0$ aux six points $x_j=0{,}02$ ; $0{,}1$ ; $0{,}25$ ; $0{,}5$ ; $0{,}75$ ; $0{,}95$. Aucun des 10 profils du front de référence ne viole ces contraintes. Sur 2\,000 profils tirés au hasard, 98~\% de ceux qui les respectent sont admissibles.
\item \textbf{Épaisseur} : $t_{\max}\ge8~\%$ est un maximum sur $x$, donc une disjonction (« il existe $x$ tel que\dots »), qui n'est pas affine. La condition suffisante $e(0{,}3)\ge0{,}08$ est affine, mais elle est violée par \textbf{7 des 10 profils du front} : ils sont minces, avec leur épaisseur maximale ailleurs. On l'a écartée.
\end{itemize}
Les contraintes de croisement entrent dans le lagrangien avec le drapeau \code{verifiee=False} : elles orientent les lectures, mais c'est le filtre géométrique exact qui juge. On ne vérifie pas une approximation quand on dispose du vrai test.

\begin{projet}
On applique ici une règle générale : \emph{dans le QUBO, ce qui oriente ; dans la vérification, ce qui garantit.} Une contrainte approchée (six points de contrôle pour une contrainte en chaque point de la corde) est utile pour orienter, dangereuse pour exclure. La condition suffisante d'épaisseur montre la limite : même comme simple orientation, une condition qui exclut 7 profils du front sur 10 détournerait la recherche.
\end{projet}

\begin{resultat}
\begin{itemize}
\item Le lagrangien garde le degré 2 même pour une contrainte quadratique, n'ajoute aucun qubit et rend le spectre à l'objectif.
\item En variables binaires, il y a un saut de dualité : on ne demande pas l'optimum au QUBO, on résout pour quelques $\lambda$, on réunit les lectures et l'on choisit en classique la meilleure lecture faisable.
\item Les multiplicateurs doivent partir petits : un terme linéaire trop fort pousse loin à l'intérieur du domaine faisable.
\item Les contraintes géométriques sont exactement affines en $z$ : elles orientent les lectures, le filtre exact juge.
\end{itemize}
\end{resultat}

\exercices
\begin{exercice}\label{ex:lag-saut}
Quatre conceptions, (marge $h$, objectif $f$) : $A=(2;\,0)$, $B=(1;\,-0{,}2)$, $C=(-1;\,-3)$, $D=(-3;\,5)$. (a) Quel est l'optimum contraint ($h\ge0$) ? (b) Écrire $L=f-\lambda h$ pour chaque conception et montrer qu'aucun $\lambda\ge0$ ne fait de $B$ le minimum. (c) Avec deux résolutions, $\lambda=0$ puis $\lambda=2$, et les deux meilleures lectures de chacune, la vérification retrouve-t-elle $B$ ? (d) Interpréter graphiquement dans le plan $(h,f)$.
\end{exercice}
\begin{exercice}\label{ex:lag-geo}
Avec $N=4$ (5 polynômes de Bernstein), $B_i(x)=\binom4i x^i(1-x)^{4-i}$. (a) Calculer les $B_i(0{,}5)$. (b) En déduire la contrainte de non-croisement en $x=0{,}5$ en fonction des coefficients $A_{{\rm ext},i}$ et $A_{{\rm int},i}$. (c) Pourquoi des coefficients d'épaisseur $A_{{\rm ext},i}-A_{{\rm int},i}$ tous positifs garantissent-ils l'absence de croisement en tout $x$ ? Pourquoi n'a-t-on pas imposé cette condition (voir le lot L6) ?
\end{exercice}

\begin{pistes}
\begin{itemize}
\item Un lagrangien \emph{augmenté} ($-\lambda h+\tfrac\rho2\min(h,0)^2$, ici sans le $\min$ pour rester quadratique) combine les deux mondes ; à tester avec un petit $\rho$.
\item Choisir $\lambda$ par dichotomie plutôt que par montée multiplicative, quand une seule contrainte domine.
\end{itemize}
\end{pistes}

\chapter{Le classifieur de convergence d'XFOIL}\label{ch:classifieur}
\iterbanner{V3}

\section{Une probabilité de succès, écrite comme un QUBO}
La convergence d'XFOIL est une contrainte cachée : on sait seulement si le calcul a abouti. La pratique de l'optimisation bayésienne contrainte est d'apprendre la probabilité de succès $p(\mathbf x)$. On choisit ici une \textbf{régression logistique} sur les mêmes 66 monômes que le modèle en $z$ :
\[ p(\mathbf z)=\sigma\big(\boldsymbol\phi(\mathbf z)^\top\boldsymbol\theta\big),\qquad\sigma(u)=\frac1{1+e^{-u}},\qquad\boldsymbol\phi(\mathbf z)=(1,z_i,z_iz_j). \]
Le \emph{log-odds} $\ln\frac p{1-p}=\boldsymbol\phi^\top\boldsymbol\theta$ est un polynôme quadratique en $z$, donc un QUBO après substitution $z=a_0+A\mathbf b$. La contrainte « $p\ge p_{\min}$ » s'écrit
\[ h_{\rm conv}(\mathbf b)=\boldsymbol\phi(\mathbf z)^\top\boldsymbol\theta-\ln\frac{p_{\min}}{1-p_{\min}}\ \ge\ 0, \]
une contrainte \emph{quadratique}. Une pénalité au carré la porterait au degré 4 ; le lagrangien la garde au degré 2. C'est la raison pour laquelle ce classifieur n'est utilisable qu'avec le lagrangien.

\paragraph{Ajustement.} On maximise la vraisemblance pénalisée $\sum_n\big[y_n\ln p_n+(1-y_n)\ln(1-p_n)\big]-\tfrac{\rho}2\|\boldsymbol\theta_{\rm sans\ constante}\|^2$, $\rho=1$, par la méthode de Newton (moindres carrés repondérés, une dizaine d'itérations). La pénalité ridge est indispensable : avec 66 coefficients et quelques dizaines d'échecs, des données séparables feraient diverger les coefficients.

\section{Ce qu'il vaut}
On rejoue les 8 exécutions de l'approche actuelle de la seconde campagne CST : pour chacun des 2\,205 profils proposés, on réapprend le classifieur sur les profils simulés \emph{avant} lui, et l'on compare sa prédiction au résultat d'XFOIL (31~\% d'échecs).
\begin{itemize}
\item \textbf{Aire sous la courbe ROC} : 0,83 (0,79 sur les 100 premières simulations, 0,85 après 200). Appris sur les 22\,860 profils admissibles du cache, il atteint 0,93 sur un jeu de test séparé : la convergence d'XFOIL est assez régulière en $z$ pour qu'un modèle quadratique la capte.
\item \textbf{Le compromis du seuil} (figure~\ref{fig:seuils}, à gauche) :
\begin{center}
\begin{tabular}{cccc}
\toprule
seuil $p_{\min}$ & échecs évités & succès sacrifiés & propositions rejetées \\
\midrule
0,1 & 31~\% & 3~\% & 12~\% \\
0,2 & 43~\% & 6~\% & 18~\% \\
0,3 & 50~\% & 8~\% & 21~\% \\
0,5 & 62~\% & 14~\% & 29~\% \\
\bottomrule
\end{tabular}
\end{center}
\end{itemize}

\begin{figure}[htbp]
\centering
% figure générée par tools/make_figures.py
\begin{tikzpicture}
\begin{groupplot}[group style={group size=2 by 1,horizontal sep=1.7cm},width=0.49\linewidth,height=5.6cm,
  xmin=0,xmax=0.55,ymin=0,ymax=105,xlabel={seuil $p_{\min}$},grid=major,grid style={cmuted!20},
  legend style={font=\scriptsize,at={(0.5,1.03)},anchor=south,legend columns=1},legend cell align=left]
\nextgroupplot[title={\small CST, rejeu de la campagne},ylabel={\%}]
\addplot[cgreen,very thick,mark=*] coordinates {(0.1,30.769) (0.2,42.899) (0.3,50.444) (0.5,62.13)};
\addlegendentry{échecs d'XFOIL évités}
\addplot[corange,very thick,mark=square*] coordinates {(0.1,3.0085) (0.2,6.2786) (0.3,8.3715) (0.5,13.865)};
\addlegendentry{succès sacrifiés}
\addplot[cink,only marks,mark=|,mark size=6pt,very thick] coordinates {(0.109,96) (0.149,96) (0.184,96) (0.356,96) (0.416,96) (0.621,96) (0.646,96) (0.724,96) (0.835,96) (0.868,96)};
\addlegendentry{$P$(succès) des 10 profils du front}
\nextgroupplot[title={\small V2, secteur S2, 20 graines}]
\addplot[cblue,very thick,mark=*] coordinates {(0.05,100) (0.1,90) (0.2,55) (0.3,30) (0.5,15)};
\addlegendentry{front exact retrouvé (\% des graines)}
\addplot[cmuted,very thick,mark=triangle*] coordinates {(0.05,59.655) (0.1,60.37) (0.2,65.448) (0.3,72.4) (0.5,75.6)};
\addlegendentry{profils simulés faisables}
\end{groupplot}
\end{tikzpicture}

\caption{Le seuil du classifieur. À gauche, sur le rejeu de la campagne CST : part des échecs d'XFOIL évités et des succès sacrifiés ; les traits verticaux (en haut) sont les probabilités de succès que donne le classifieur appris sur tout le cache aux 10 profils du front de référence. À droite, sur le problème V2 (S2) : part des graines qui retrouvent le front exact, et part des profils simulés faisables.}\label{fig:seuils}
\end{figure}

\section{Le front est au bord du domaine d'XFOIL}
Le chiffre qui compte est ailleurs. Le classifieur appris sur \emph{tout} le cache donne aux 10 profils du front de référence des probabilités de succès de 0,11 ; 0,15 ; 0,18 ; 0,36 ; 0,42 ; 0,62 ; 0,65 ; 0,72 ; 0,83 et 0,87. Trois sont sous 0,2. Sur le problème V2, le profil extrême du front S2, $(7,0,2,0)$, a 4 voisins sur 5 en échec, et le classifieur lui donne 0,2. Les profils du front sont minces et peu épais au bord d'attaque : c'est là que la traînée et la SER sont faibles, et aussi là qu'XFOIL décroche.

Conséquence directe sur le problème V2 (figure~\ref{fig:seuils}, à droite) : avec $p_{\min}=0{,}5$, la boucle propose 76~\% de faisables au lieu de 60~\%, et atteint 99~\% de l'hypervolume plus tôt (62,5 simulations au lieu de 75), mais elle ne retrouve le front exact que pour 15~\% des graines. Avec $p_{\min}=0{,}05$, elle se comporte comme le lagrangien seul.

\begin{piege}[ : un classifieur de contrainte cachée n'est pas une barrière]
Un seuil élevé rend la boucle plus économe et plus aveugle : il exclut les profils voisins des échecs, et le front y est. C'est la même leçon que le réseau de neurones de la V2, sous une autre forme. La pratique industrielle (critère multiplié par la probabilité de faisabilité) n'exclut jamais : elle \emph{pondère}. Pour un classifieur utilisé comme contrainte, le seuil doit rester bas ($p_{\min}=0{,}1$ dans la campagne du chapitre~\ref{ch:campagne}).
\end{piege}

\begin{projet}
Que le front soit au bord du domaine de validité d'XFOIL est un résultat physique, pas seulement numérique. Un profil voisin d'échecs de convergence est proche du décollement, et sa traînée calculée est moins fiable. C'est un argument pour le lot L3 (robustesse physique) : vérifier les profils du front par un calcul plus fidèle, ou les pénaliser par leur distance au domaine d'échec.
\end{projet}

\begin{resultat}
\begin{itemize}
\item Une régression logistique quadratique en $z$ donne une contrainte de convergence quadratique, utilisable avec le lagrangien seulement. Elle prédit bien les échecs (AUC 0,83 en cours de boucle).
\item Le front est collé au domaine où XFOIL échoue : un seuil élevé exclut ses extrémités. Le classifieur doit orienter (seuil bas), pas exclure.
\end{itemize}
\end{resultat}

\exercices
\begin{exercice}\label{ex:clf-budget}
Rejeu CST : sur 300 propositions, 31~\% échouent. (a) Avec $p_{\min}=0{,}2$, combien d'échecs évite-t-on et combien de succès sacrifie-t-on ? (b) Bilan en profils faisables, si chaque proposition rejetée est remplacée par une proposition au même taux de succès que les autres ? (c) Quel risque ce bilan ne compte-t-il pas ?
\end{exercice}
\begin{exercice}\label{ex:clf-seuil}
$\ln\frac{p}{1-p}$ pour $p_{\min}=0{,}1$ ; $0{,}2$ ; $0{,}5$. Un profil a un log-odds prédit de $-1{,}8$ : pour quels seuils est-il accepté ? Que vaut sa probabilité de succès prédite ?
\end{exercice}

\begin{pistes}
\begin{itemize}
\item Pondérer au lieu d'exclure : choisir parmi les lectures vérifiées en maximisant (amélioration prédite) $\times\ p$, comme en optimisation bayésienne contrainte.
\item Apprendre le classifieur avec les échecs de \emph{toutes} les exécutions (le cache) : la convergence d'XFOIL ne dépend pas de la méthode d'optimisation, et l'AUC passe de 0,83 à 0,93.
\end{itemize}
\end{pistes}

\chapter{Ce que l'on mesure sans XFOIL}\label{ch:hors-ligne}
\iterbanner{V3}

Trois études ne demandent aucune simulation nouvelle : le problème V2, dont toutes les simulations sont tabulées ; le rejeu de la seconde campagne CST à partir du cache ; et QAOA sur de petites instances. Toutes utilisent le modèle en $z$ pour l'objectif.

\section{Le problème V2 : la boucle complète, 20 graines}
Même protocole que les lots L1 et L6 : 10 profils initiaux identiques par graine, 5 poids, budget 250, solveur exact (\code{contraintes/experiences\_v2.py}). On compte les simulations pour retrouver le front exact (251 si jamais retrouvé) et pour atteindre 99~\% de son hypervolume.
\begin{center}
\renewcommand{\arraystretch}{1.2}
\small
\begin{tabularx}{\linewidth}{>{\raggedright\arraybackslash}X c c c c c}
\toprule
& front complet & succès & HV 99~\% & faisables & vs bits d'écart \\
\midrule
\multicolumn{6}{l}{\textit{S1 (front exact : 4 profils)}}\\
bits d'écart (actuel) & 75 (60--96) & 100~\% & 62,5 & 53~\% & --- \\
lagrangien vérifié & 77,5 (65--131) & 100~\% & 57,5 & 58~\% & 9 mieux, 10 pire, $p=0{,}50$ \\
lagrangien + classifieur ($p_{\min}=0{,}5$) & 85 (60--106) & 100~\% & 60 & 64~\% & 7 / 11, $p=0{,}38$ \\
déséquilibrée & 120 (70--146) & 100~\% & 65 & 58~\% & 5 / 14, $p=0{,}03$ \\
\midrule
\multicolumn{6}{l}{\textit{S2 (front exact : 10 profils)}}\\
bits d'écart (actuel) & 142,5 (116--191) & 100~\% & 75 & 61~\% & --- \\
lagrangien vérifié & \textbf{102,5} (90--116) & 100~\% & 75 & 60~\% & \textbf{17 / 3}, $p=0{,}0002$ \\
lagrangien + classifieur ($p_{\min}=0{,}5$) & --- & 15~\% & 62,5 & 76~\% & 2 / 18, $p=0{,}002$ \\
déséquilibrée & --- & 0~\% & 72,5 & 60~\% & 0 / 20, $p<0{,}001$ \\
\bottomrule
\end{tabularx}
\end{center}
(Médianes, entre parenthèses les quartiles ; test de Wilcoxon apparié sur le nombre de simulations pour le front complet.)

\paragraph{Lecture.} Avec le solveur exact, les lectures sont les 600 meilleures chaînes : la vérification choisit alors l'optimum \emph{exact} du problème contraint par les modèles, quels que soient les $\lambda$. La ligne « lagrangien vérifié » mesure donc ce que coûte la pénalité à bits d'écart, à modèles identiques. En S1, rien ; en S2, dont le front longe la contrainte de manœuvre (lot L1), elle fait perdre 40 simulations en médiane : les petites violations rentables du chapitre~\ref{ch:penalites} envoient la boucle sur des profils que le modèle lui-même jugeait infaisables. La pénalité déséquilibrée, qui vise une marge de 0,05, ne retrouve jamais le front S2 : ses profils sont collés à la frontière. Le classifieur avec un seuil de 0,5 coupe les extrémités du front (chapitre~\ref{ch:classifieur}).

\paragraph{Une question du lot L1 résolue.} Le lot L1 avait trouvé qu'un « contrôle sans QUBO » (mêmes modèles, minimum de la prédiction cherché directement) retrouvait le front S2 en 95 simulations, contre 142,5 pour la boucle QUBO, sans pouvoir isoler la cause : variable d'écart discrète, régularisation ou filtre d'épaisseur. Le lagrangien vérifié, qui garde le QUBO mais retire la variable d'écart, fait 102,5. L'essentiel de l'écart venait donc de la pénalité à bits d'écart, pas de la forme QUBO.

\section{Le problème CST : rejeu de la seconde campagne}
On rejoue les 8 exécutions de l'approche actuelle (\code{contraintes/etude\_cst.py}) : pour chaque profil proposé, on réajuste les modèles sur les profils simulés avant lui, comme la boucle à cette itération.

\paragraph{D'où viennent les $C_{L,\max}$ trop faibles ?} Sur 1\,609 profils proposés où XFOIL a convergé, 678 (42~\%) ont $C_{L,\max}<1{,}2$. Pour \textbf{317} d'entre eux, le modèle affine prédisait \emph{déjà} moins de 1,2 : la contrainte était connue du modèle, et c'est la formulation qui ne l'a pas fait respecter. Pour les 361 autres, le modèle s'est trompé (erreur absolue médiane 0,15, biais $+0{,}02$). La moitié des $C_{L,\max}$ trop faibles est donc un défaut de formulation, l'autre moitié un défaut de modèle.

\paragraph{Les formulations sur les mêmes QUBO.} À 60, 150 et 300 simulations des 8 exécutions, pour les 5 poids (120 QUBO), on construit le QUBO de chaque formulation, on le résout par recuit simulé (128 lectures), et l'on retient le candidat comme le ferait la boucle.
\begin{center}
\renewcommand{\arraystretch}{1.2}
\small
\begin{tabular}{lcccc}
\toprule
& qubits & part de l'objectif & candidat prédit faisable & $\hat C_{L,\max}$ du candidat \\
\midrule
bits d'écart (actuel) & 33 & 0,8~\% & 61~\% & 1,20 \\
déséquilibrée & 30 & 2,4~\% & 97~\% & 1,25 \\
lagrangien vérifié & 30 & $\approx100$~\% & 92~\% & 1,23 à 1,34 \\
\bottomrule
\end{tabular}
\end{center}
Comparés deux à deux sur les QUBO où les deux candidats sont prédits faisables, le lagrangien retient un meilleur objectif que les bits d'écart dans 63~\% des cas (68 QUBO), et que la pénalité déséquilibrée dans 75~\% (110 QUBO) ; les écarts d'objectif sont faibles (0,015 en médiane, pour une étendue de l'ordre de 1). La différence est donc d'abord la \textbf{faisabilité} : les bits d'écart rendent un candidat que leur propre modèle juge infaisable 4 fois sur 10.

\section{QAOA sur de petites instances}
Mêmes QUBO de début de boucle que le cours IA et QUBO (V2, S2, 30 profils, 3 graines, 3 poids ; \code{contraintes/qaoa\_formulations.py}). Pour comparer des QUBO de tailles différentes, on mesure la probabilité, sous QAOA ($p=2$), d'obtenir la \emph{conception} optimale du problème contraint (bits auxiliaires marginalisés), rapportée au tirage uniforme d'une conception ($1/1024$).
\begin{center}
\renewcommand{\arraystretch}{1.2}
\small
\begin{tabular}{lcccc}
\toprule
& qubits & part de l'objectif & gain sur l'optimum contraint & échantillons faisables \\
\midrule
bits d'écart (actuel) & 14 & 1,7~\% & 1,4 (1,0 à 1,8) & 80~\% \\
déséquilibrée & 11 & 3,8~\% & 1,9 (0,0 à 3,3) & 64~\% \\
lagrangien & 10 & 100~\% & \textbf{10,3} (0,0 à 72,9) & 37~\% \\
\bottomrule
\end{tabular}
\end{center}
\paragraph{Lecture.} Le lagrangien multiplie par 7 en médiane la probabilité que QAOA trouve l'optimum contraint, mais avec une forte dispersion. La vérification explique les deux. Dans 8 instances sur 9, \emph{aucun} $\lambda$ ne fait de l'optimum contraint l'état fondamental du lagrangien : c'est le saut de dualité, aggravé ici par l'épaisseur, laissée hors du QUBO. Mais dans 7 cas sur 9 l'optimum est parmi les 17 états les plus bas du spectre (sur 1\,024), et QAOA, qui répartit sa probabilité sur les états de basse énergie, le trouve souvent. Dans les deux cas où $\lambda$ est grand (3,2), l'optimum est au rang 196 ou 320, et QAOA ne le voit pas. Les échantillons sont moins souvent faisables (37~\%), ce qui ne coûte rien : la vérification les trie. Avec le solveur exact, la vérification retrouve l'optimum contraint dans les 9 cas.

\begin{resultat}
\begin{itemize}
\item V2 : le lagrangien vérifié retrouve le front S2 en 102 simulations au lieu de 142 (17 graines sur 20) ; c'est le coût de la pénalité à bits d'écart, à modèles identiques.
\item CST : la moitié des $C_{L,\max}$ trop faibles vient de la formulation ; les bits d'écart rendent un candidat que leur propre modèle juge infaisable 4 fois sur 10, le lagrangien moins d'une fois sur 10.
\item QAOA : sans pénalité, l'objectif retrouve le spectre et QAOA trouve l'optimum contraint 7 fois plus souvent en médiane, à condition que $\lambda$ reste petit.
\end{itemize}
\end{resultat}

\exercices
\begin{exercice}\label{ex:hl-exact}
Pourquoi, avec le solveur exact et la vérification, les multiplicateurs n'ont-ils aucune influence sur le profil retenu dans le problème V2 ? Combien de chaînes faudrait-il garder pour que ce soit encore vrai à 30 bits, et pourquoi est-ce impossible ?
\end{exercice}
\begin{exercice}\label{ex:hl-gain}
Dans le QUBO à bits d'écart (14 qubits), la conception optimale peut être associée à 16 valeurs des bits auxiliaires. (a) Pourquoi marginaliser sur ces bits pour comparer à un QUBO de 10 qubits ? (b) Un gain de 1 signifie-t-il que QAOA ne fait rien ? (c) Pourquoi le gain mesuré dans le cours IA et QUBO (4,5) diffère-t-il de celui-ci (1,4) pour la même formulation ?
\end{exercice}

\begin{pistes}
\begin{itemize}
\item Rejouer de même les exécutions de la machine de factorisation et du modèle bayésien : les défauts de formulation devraient être les mêmes, puisque la contrainte est la même.
\item Mesurer le gain de QAOA sur la meilleure \emph{lecture vérifiée} parmi $k$ tirs, plutôt que sur la probabilité de l'optimum : c'est la quantité utile dans la boucle.
\end{itemize}
\end{pistes}

\chapter{L'interface des solveurs industriels}\label{ch:cqm}
\iterbanner{V3}

\section{Garder les contraintes séparées}
Jusqu'ici, la boucle livrait à son solveur une seule matrice $Q$ où pénalités et objectif sont mélangés. Les solveurs industriels attendent autre chose : un \textbf{modèle quadratique contraint} (en anglais \emph{constrained quadratic model}, CQM),
\[ \min_{\mathbf b\in\{0,1\}^n}\ \mathbf b^\top Q\,\mathbf b+c\quad\text{sous}\quad\mathbf b^\top H_k\mathbf b+h_{0,k}\ \ge\ 0,\quad k=1,\dots,K. \]
C'est ce que reçoivent les solveurs classiques en variables binaires (CPLEX, Gurobi, SCIP) et le solveur hybride contraint de D-Wave, qui accepte directement contraintes d'égalité et d'inégalité. C'est \emph{lui} qui choisit comment traiter les contraintes : pénalités réglées automatiquement, décomposition, prétraitement classique qui réduit le problème avant la partie quantique. Les retours d'expérience industriels (planification de la production de véhicules d'essai, par exemple) soulignent l'avantage principal : on n'a plus à régler soi-même les poids de pénalité, réglage qui devient très difficile quand le modèle grossit.

\section{Notre export}
\code{contraintes/export\_cqm.py} écrit le problème d'une itération au \textbf{format LP}, un format texte lu par tous ces outils (côté D-Wave, par \code{dimod.lp.load}), et en JSON. L'exemple fourni est une itération réelle de la seconde campagne CST (graine 0, 150 profils simulés, $w=0{,}5$) :
\begin{itemize}
\item 30 variables binaires, nommées \code{a0\_b0} à \code{a9\_b2} (coefficient, bit) ;
\item l'objectif pondéré normalisé, quadratique ;
\item 8 contraintes : $\hat C_{L,\max}\ge1{,}2$ (linéaire), non-croisement aux 6 points de contrôle (linéaires), convergence d'XFOIL $p\ge0{,}1$ (quadratique).
\end{itemize}

\begin{codebox}[ : début de \code{Code/resultats/cqm\_exemple.lp}]
\begin{lstlisting}[language={}]
\ Boucle QUBO, problème CST 30 bits, graine 0, 150 profils simulés, w = 0,5 (lot L7)
Minimize
 obj: - 0.0927 a0_b0 - 0.148 a0_b1 ... + [ + 0.148 a0_b0 * a0_b1 + ... ] / 2
Subject To
 c0_CLmax_min: ... >= ...
 c1_croisement_x_0_02: + 0.0454 a0_b0 + 0.0907 a0_b1 + ... >= ...
 ...
 c7_convergence_xfoil: ... + [ ... ] >= ...
Binary
 a0_b0 a0_b1 a0_b2 a1_b0 ...
End
\end{lstlisting}
\end{codebox}
\paragraph{Validation.} Relu par \code{dimod} 0.12 (installé seulement pour ce contrôle, ce n'est pas une dépendance du projet), le fichier redonne 30 variables et 8 contraintes ; sur 200 chaînes tirées au hasard, l'objectif et les 8 contraintes coïncident avec ceux de la boucle à $10^{-9}$ près.

\begin{piege}[ : les conventions du format LP]
Dans l'objectif, la partie quadratique entre crochets est \emph{divisée par 2} : il faut y écrire $2Q_{ij}$. Dans les contraintes, elle ne l'est pas. La constante de l'objectif n'est pas écrite (certains lecteurs la refusent) ; elle ne change pas l'optimum et figure en commentaire. Une erreur de facteur 2 sur les couplages donnerait un fichier lisible mais faux : d'où la vérification sur des chaînes tirées au hasard.
\end{piege}

\section{Ce que cela change pour l'équipe}
\begin{itemize}
\item Le problème d'une itération devient un \emph{fichier} que l'on peut confier à n'importe quel solveur, quantique, hybride ou classique, et comparer à budget égal, sans toucher à la boucle.
\item La comparaison la plus utile pour le jury : résoudre le même CQM par le recuit simulé de la boucle, par un solveur exact classique, et par un solveur hybride quantique. Si les trois rendent le même candidat, le choix du solveur n'est pas le goulot ; s'ils diffèrent, on sait où chercher.
\item Pour un industriel, c'est le format attendu : on ne lui livre pas une matrice pré-pénalisée dont il ne peut pas relire les contraintes.
\end{itemize}

\begin{resultat}
\begin{itemize}
\item Les solveurs industriels reçoivent objectif et contraintes séparés et règlent eux-mêmes le traitement des contraintes.
\item L'export LP de la boucle est validé avec \code{dimod} ; il rend la chaîne compatible avec les solveurs classiques et hybrides sans modifier la boucle.
\end{itemize}
\end{resultat}

\exercices
\begin{exercice}\label{ex:cqm-ecrire}
Écrire au format LP le problème : minimiser $b_0-2b_1+3b_0b_1$ sous $b_0+b_1\ge1$ et $b_0b_1-0{,}5\,b_0\ge-0{,}2$, $b_0,b_1$ binaires. Résoudre à la main et vérifier.
\end{exercice}

\begin{pistes}
\begin{itemize}
\item Résoudre le CQM exemple avec un solveur exact (SCIP, libre) pour connaître l'optimum contraint du modèle à 30 bits, et mesurer l'écart du recuit simulé vérifié.
\item Le soumettre au solveur hybride contraint de D-Wave si l'équipe obtient un accès (programme académique), et comparer les candidats.
\end{itemize}
\end{pistes}

\chapter{La campagne avec XFOIL}\label{ch:campagne}
\iterbanner{V3}

\section{Protocole}
Exactement celui de la seconde campagne CST du lot L6 (\code{contraintes/experiences\_cst.py}) : filtre géométrique avant simulation, traînée de plus de 200 counts comptée comme un échec, budget de 300 simulations, 8 graines, \emph{mêmes profils initiaux par graine}, ce qui permet une comparaison appariée. Les méthodes de la seconde campagne ne sont pas relancées : leurs exécutions sont relues, et le front de référence est recalculé sur l'union des deux campagnes. Deux méthodes nouvelles, toutes deux avec le modèle en $z$ et le recuit simulé, à 30 qubits :
\begin{itemize}
\item \code{quadz\_lagrange} : $C_{L,\max}$ et les six contraintes de non-croisement en termes lagrangiens, 6 résolutions par QUBO, lectures vérifiées ;
\item \code{quadz\_lagrange\_clf} : la même chose, plus le classifieur de convergence d'XFOIL avec $p_{\min}=0{,}1$.
\end{itemize}
La pénalité déséquilibrée, moins bonne que le lagrangien sur les deux études hors ligne, n'a pas été relancée.

\section{Résultats}
Hypervolume des profils faisables simulés, rapporté à celui du front de référence (11 profils, sur l'union des 20\,602 profils simulés par les deux campagnes), médiane sur 8 graines. « vs actuel » : nombre de graines où la méthode fait mieux que l'approche actuelle à 300 simulations, et $p$ du test de Wilcoxon apparié. À 100 simulations, le classifieur fait mieux sur 5 graines ($p=0{,}84$).
\begin{center}
\renewcommand{\arraystretch}{1.2}
\setlength{\tabcolsep}{4pt}
\small
\begin{tabular}{lcccccc}
\toprule
& à 100 & à 300 & vs actuel (300) & faisables & échecs XFOIL & $C_{L,\max}$ faible \\
\midrule
bits d'écart (actuel) & 0,768 & \textbf{0,882} & --- & 43~\% & 28~\% & 29~\% \\
lagrangien vérifié & 0,768 & 0,839 & 2/8, $p=0{,}38$ & 37~\% & 40~\% & 22~\% \\
lagrangien + classifieur & \textbf{0,815} & 0,840 & 1/8, $p=0{,}15$ & 41~\% & 34~\% & 25~\% \\
\midrule
factorisation (L6) & 0,833 & 0,895 & 6/8, $p=0{,}38$ & 47~\% & 25~\% & 28~\% \\
NSGA-II & 0,718 & 0,813 & 2/8, $p=0{,}04$ & 72~\% & 13~\% & 15~\% \\
hasard & 0,683 & 0,772 & 1/8, $p=0{,}02$ & 76~\% & 6~\% & 17~\% \\
\bottomrule
\end{tabular}
\end{center}

\begin{figure}[htbp]
\centering
% figure générée par tools/make_figures.py
\begin{tikzpicture}
\begin{axis}[width=0.92\linewidth,height=7cm,xmin=0,xmax=300,ymin=0.3,ymax=1,xlabel={simulations},
  ylabel={$HV/HV_{\rm ref}$ (médiane sur 8 graines)},grid=major,grid style={cmuted!20},legend pos=south east,
  legend cell align=left,legend style={font=\footnotesize}]
\addplot[cmuted,thick,dashed,mark size=1.4pt] coordinates {(20,0.61544) (40,0.64622) (60,0.65565) (80,0.66783) (100,0.68305) (120,0.6884) (140,0.72706) (160,0.72737) (180,0.72737) (200,0.72913) (220,0.76737) (240,0.77032) (260,0.77032) (280,0.77082) (300,0.77155)};
\addlegendentry{hasard}
\addplot[cink,thick,dotted,mark size=1.4pt] coordinates {(20,0.61096) (40,0.62004) (60,0.6257) (80,0.65257) (100,0.71827) (120,0.72088) (140,0.74279) (160,0.76578) (180,0.77044) (200,0.77275) (220,0.78817) (240,0.78883) (260,0.789) (280,0.81253) (300,0.81278)};
\addlegendentry{NSGA-II}
\addplot[cblue,very thick,mark=*,mark size=1.4pt] coordinates {(20,0.61204) (40,0.69689) (60,0.72471) (80,0.76812) (100,0.76812) (120,0.76812) (140,0.79083) (160,0.80248) (180,0.82308) (200,0.82499) (220,0.82879) (240,0.84877) (260,0.85732) (280,0.86329) (300,0.88181)};
\addlegendentry{bits d'écart (actuel)}
\addplot[corange,thick,mark=square*,mark size=1.4pt] coordinates {(20,0.63019) (40,0.6879) (60,0.76189) (80,0.8174) (100,0.83258) (120,0.85689) (140,0.86559) (160,0.86807) (180,0.87256) (200,0.87577) (220,0.87915) (240,0.8796) (260,0.88328) (280,0.89162) (300,0.89483)};
\addlegendentry{machine de factorisation (bits d'écart)}
\addplot[cgreen,very thick,mark=triangle*,mark size=1.4pt] coordinates {(20,0.58278) (40,0.62809) (60,0.708) (80,0.75695) (100,0.7682) (120,0.79466) (140,0.79466) (160,0.79466) (180,0.79466) (200,0.79986) (220,0.83212) (240,0.83255) (260,0.83795) (280,0.83895) (300,0.83944)};
\addlegendentry{lagrangien vérifié}
\addplot[cviolet,very thick,mark=diamond*,mark size=1.4pt] coordinates {(20,0.58278) (40,0.69995) (60,0.75154) (80,0.80242) (100,0.81519) (120,0.82286) (140,0.82286) (160,0.82286) (180,0.82286) (200,0.82286) (220,0.82644) (240,0.82822) (260,0.836) (280,0.83936) (300,0.83991)};
\addlegendentry{lagrangien + classifieur}
\end{axis}
\end{tikzpicture}

\caption{Campagne CST : hypervolume relatif en fonction du nombre de simulations (médiane sur 8 graines). Les deux méthodes lagrangiennes (vert, violet) partent comme l'approche actuelle, puis plafonnent.}\label{fig:campagne-hv}
\end{figure}

\paragraph{1. Les études hors ligne ne se confirment pas.} Les deux méthodes lagrangiennes finissent \emph{sous} l'approche actuelle (0,84 contre 0,88), sans que l'écart soit significatif avec 8 graines. Le classifieur aide en début de budget (0,815 à 100 simulations, devant l'approche actuelle sur 5 graines), puis plafonne aussi.

\paragraph{2. Ce qui était prévu se produit.} Les $C_{L,\max}$ trop faibles baissent, de 29~\% à 22~\% : la formulation fait mieux respecter la contrainte du modèle, comme l'annonçait le rejeu (chapitre~\ref{ch:hors-ligne}). Le classifieur fait baisser les échecs d'XFOIL de 40~\% à 34~\% par rapport au lagrangien seul.

\paragraph{3. Ce qui n'était pas prévu : les échecs d'XFOIL montent.} De 28~\% à 40~\% avec le lagrangien seul. La faisabilité totale baisse donc (37~\% contre 43~\%). On a cherché pourquoi dans les profils proposés : la part des coefficients placés à un niveau \emph{extrême} de la grille (0 ou 7).
\begin{center}
\begin{tabular}{lc}
\toprule
& coefficients aux niveaux extrêmes \\
\midrule
hasard & 25~\% \\
machine de factorisation (régularisée) & 30~\% \\
bits d'écart (actuel) & 56~\% \\
lagrangien + classifieur & 62~\% \\
lagrangien vérifié & \textbf{66~\%} \\
\bottomrule
\end{tabular}
\end{center}
Pour les trois méthodes à modèle en $z$, les profils dont au moins 6 coefficients sur 10 sont extrêmes échouent dans XFOIL 40 à 48~\% du temps, contre 16 à 21~\% pour les autres.

\paragraph{4. Le mécanisme.} Le modèle en $z$ est une parabole à 66 coefficients, sans régularisation : son minimum est souvent sur un bord de la grille, là où il extrapole. Avec les bits d'écart, la pénalité occupait 99~\% du spectre ; le recuit, réglé sur l'objectif, l'optimisait mal, et ses lectures restaient plus centrées. \emph{La pénalité régularisait par accident.} Le lagrangien rend tout le spectre à l'objectif : le recuit trouve beaucoup mieux le minimum du modèle, c'est-à-dire ses erreurs d'extrapolation dans les coins, où XFOIL décroche. La machine de factorisation, régularisée, ne va aux extrêmes que 30~\% du temps : c'est elle qu'il faudrait associer au lagrangien.

\paragraph{5. Pourquoi le rejeu ne l'a pas vu.} Il mesurait la faisabilité \emph{prédite} et l'objectif \emph{prédit} du candidat, sur des états produits par l'approche actuelle. Or le défaut du lagrangien n'est pas dans le modèle de contrainte : il est dans la confiance excessive qu'il accorde au modèle d'objectif, ce qu'aucune prédiction ne peut révéler. Seule la simulation le montre.

\begin{piege}[ : mieux résoudre un mauvais modèle]
Résoudre plus exactement le problème du modèle n'aide que si le modèle est bon là où l'on va. Un solveur meilleur, quantique ou non, appliqué à un modèle qui extrapole, trouve plus sûrement ses erreurs. C'est l'argument classique des \emph{régions de confiance} en optimisation par modèle de substitution (TuRBO, par exemple) : limiter la recherche là où le modèle a des données.
\end{piege}

\begin{resultat}
\begin{itemize}
\item Sur le problème CST, avec XFOIL, le lagrangien vérifié ne bat pas les bits d'écart : 0,84 contre 0,88 d'hypervolume relatif à 300 simulations (non significatif).
\item Il respecte mieux la contrainte de $C_{L,\max}$, mais il pousse le modèle en $z$ dans ses coins, où XFOIL échoue : la pénalité à bits d'écart régularisait par accident.
\item Le classifieur réduit les échecs et aide en début de budget, sans changer le résultat final.
\item Une étude hors ligne sur des prédictions ne peut pas voir un défaut qui tient à la confiance accordée au modèle d'objectif.
\end{itemize}
\end{resultat}

\exercices
\begin{exercice}\label{ex:camp-signe}
Avec 8 graines et un test de signe bilatéral, à partir de combien de graines favorables sur 8 un écart devient-il significatif au seuil de 5~\% ? Que peut-on conclure d'un résultat « 2 graines sur 8 » ?
\end{exercice}
\begin{exercice}\label{ex:camp-coins}
Une parabole en une variable, $\hat y(z)=a z^2+bz+c$ sur $z\in[-1,1]$, ajustée sur des points tous situés dans $[-0{,}5;\,0{,}5]$. (a) Si $a<0$, où est son minimum sur $[-1,1]$ ? (b) Que vaut l'incertitude de la prédiction en $z=\pm1$ comparée à celle en $z=0$ ? (c) Expliquer pourquoi un solveur qui trouve exactement ce minimum est plus exposé qu'un solveur bruité.
\end{exercice}

\begin{pistes}
\begin{itemize}
\item Campagne \code{fm\_lagrange\_clf} (8 graines, une quarantaine de minutes) : la machine de factorisation va peu aux extrêmes, c'est l'association naturelle.
\item Une région de confiance dans le QUBO : la distance de Hamming à un profil de référence est \emph{linéaire} en bits, donc une contrainte « au plus $k$ bits changés » entre directement dans le lagrangien.
\item Régulariser le modèle en $z$ (ridge plus fort) avant de lui rendre le spectre.
\end{itemize}
\end{pistes}


\chapter{Bilan}\label{ch:bilan}
\iterbanner{V3}
\textit{Bilan à venir.}


\appendix
\chapter{Corrigés des exercices}\label{annexe:corriges}
\iterbanner{V3}

\section*{Chapitre~\ref{ch:constat} — Le constat}
\textbf{Ex.~\ref{ex:constat-budget}.} (a) $300\times0{,}43=129$ profils faisables. (b) Les pertes passent de 57~\% à 28,5~\% : $300\times0{,}715\approx215$ faisables, ce que l'approche actuelle n'obtiendrait qu'avec $215/0{,}43\approx500$ simulations. (c) Le front est au bord des contraintes : une méthode qui le cherche propose forcément des profils de part et d'autre de la frontière ; les modèles appris sur quelques centaines de points se trompent près de la frontière ; et la convergence d'XFOIL n'est connue qu'après simulation. Le hasard fait 76~\% parce qu'il ne cherche pas le bord.

\section*{Chapitre~\ref{ch:taxonomie} — Classer les contraintes}
\textbf{Ex.~\ref{ex:taxo-classer}.} (a) Envergure : connue, a priori, quantifiable, relâchable (le code calcule au-delà) ; on l'impose par les bornes de la grille, sans pénalité. (b) Mailleur qui échoue : cachée, non quantifiable, non relâchable ; mais le mailleur est rapide, donc on le lance \emph{avant} le solveur, comme un filtre, et l'on apprend un classifieur pour orienter. (c) Flutter : connue, issue de la simulation, quantifiable (la marge), relâchable ; un modèle de substitution à part, comme $C_{L,\max}$. (d) Solution instationnaire : connue mais non quantifiable si l'on n'a qu'un indicateur, quantifiable si l'on mesure l'amplitude des oscillations ; issue de la simulation ; les sorties existent mais n'ont pas de sens (non relâchable en pratique) : un classifieur, ou un modèle de l'amplitude.

\textbf{Ex.~\ref{ex:taxo-filtre}.} Budget : une pénalité finie, sur un modèle approché, n'empêche pas de proposer un profil qui la viole ; chaque proposition de ce type est une simulation perdue, et comme la contrainte n'est pas relâchable, elle ne rend même pas d'information. Modèle : les profils filtrés n'entrent jamais dans les données ; avec une pénalité, la boucle simule ces profils, ou au contraire n'apprend rien près de la frontière, et le modèle extrapole. Le filtre est exact et gratuit, la pénalité est approchée et coûteuse.

\section*{Chapitre~\ref{ch:penalites} — Pénalités}
\textbf{Ex.~\ref{ex:pen-rentable}.} (a) Le profil faisable vaut $-0{,}80$. Le voisin vaut $-1{,}00+75\times0{,}04^2=-1{,}00+0{,}12=-0{,}88$ : c'est lui, infaisable, l'optimum du QUBO. (b) Il faut $P\times0{,}0016>0{,}20$, soit $P>125$. (c) Pour $h=-0{,}5$, la pénalité vaut $125\times0{,}25\approx31$, pour un objectif d'étendue 1 : l'objectif n'occupe plus que 3~\% environ du spectre. Rendre la pénalité sûre pour les petites violations la rend écrasante pour les grandes.

\textbf{Ex.~\ref{ex:pen-deseq}.} (a) $75(h^2-0{,}1h)$ : $0{,}56$ en $h=-0{,}05$ ; $0$ en $h=0$ ; $-0{,}19$ en $h=0{,}05$ ; $4{,}5$ en $h=0{,}3$. (b) Pour $h>2h^\star=0{,}1$ : tout profil dont la marge dépasse 0,1 est pénalisé. (c) Avec $h^\star=0{,}15$, la pénalité s'annule en $h=0{,}3$ quel que soit $P$ ; avec $P=30$, elle vaut $0{,}53$ en $h=-0{,}05$ et $-0{,}68$ en $h=0{,}15$. On ne pénalise plus les profils de marge $0{,}3$, mais on attire maintenant les candidats vers une marge de 0,15, loin du front qui longe la contrainte. On ne peut pas avoir les deux avec une parabole.

\textbf{Ex.~\ref{ex:pen-degre}.} Avec $z_i=2b_i-1$ : $h=4b_1b_2-2b_1-2b_2+0{,}8$. Dans $h^2$, les produits comme $(b_1b_2)^2$ ou $b_1^2b_2$ se réduisent à $b_1b_2$, puisque $b^2=b$ : avec deux bits seulement, tout polynôme est au plus de degré 2. Avec 3 bits par variable, $z_1=\sum_jc_jb_{1j}+c_0$ et $z_1z_2$ contient les produits $b_{1j}b_{2k}$ ; son carré contient $b_{1j}b_{1j'}b_{2k}b_{2k'}$ avec $j\ne j'$ et $k\ne k'$, de degré 4, que rien ne réduit.

\section*{Chapitre~\ref{ch:lagrangien} — Lagrangien}
\textbf{Ex.~\ref{ex:lag-saut}.} (a) Faisables : $A$ ($f=0$) et $B$ ($f=-0{,}2$) : l'optimum contraint est $B$. (b) $L_A=-2\lambda$, $L_B=-0{,}2-\lambda$, $L_C=-3+\lambda$, $L_D=5+3\lambda$. $L_B\le L_A$ exige $\lambda\le0{,}2$ ; $L_B\le L_C$ exige $\lambda\ge1{,}4$. C'est incompatible : $B$ n'est jamais le minimum. (c) $\lambda=0$ : $L=(0;\,-0{,}2;\,-3;\,5)$, les deux meilleures sont $C$ et $B$. $\lambda=2$ : $L=(-4;\,-2{,}2;\,-1;\,11)$, les deux meilleures sont $A$ et $B$. La réunion $\{A,B,C\}$ contient $B$, que la vérification retient (faisable, meilleur objectif). (d) Minimiser $f-\lambda h$ revient à pousser vers le haut une droite de pente $\lambda$ jusqu'à toucher un point : seuls les points de l'enveloppe convexe inférieure peuvent être touchés. $B=(1;\,-0{,}2)$ est au-dessus du segment $AC$ (qui vaut $-1$ en $h=1$) : il est invisible pour tout $\lambda$, mais il fait partie des bonnes lectures.

\textbf{Ex.~\ref{ex:lag-geo}.} (a) $B_i(0{,}5)=\binom4i/16$ : $\tfrac1{16},\tfrac4{16},\tfrac6{16},\tfrac4{16},\tfrac1{16}$. (b) $\Delta_0+4\Delta_1+6\Delta_2+4\Delta_3+\Delta_4\ge0$ avec $\Delta_i=A_{{\rm ext},i}-A_{{\rm int},i}$, chaque $\Delta_i$ étant affine en $z_i$ et $z_{i+5}$. (c) Les $B_i$ et $C$ sont positifs sur $[0,1]$ : une somme de termes positifs est positive. On ne l'a pas imposé parce que c'est une condition suffisante trop sévère : le lot L6 a montré que 3 des 10 profils du front de référence ont un $\Delta_i$ négatif tout en étant valides.

\section*{Chapitre~\ref{ch:classifieur} — Classifieur}
\textbf{Ex.~\ref{ex:clf-budget}.} (a) 93 échecs et 207 succès ; on évite $0{,}43\times93\approx40$ échecs et l'on sacrifie $0{,}063\times207\approx13$ succès. (b) 53 propositions rejetées ; remplacées par des propositions qui passent le classifieur, dont le taux de succès est $(207-13)/(300-53)\approx0{,}79$, elles donnent environ 42 succès : $194+42\approx236$ faisables au lieu de 207, soit 14~\% de plus. (c) Les succès sacrifiés ne se valent pas : si ce sont des profils du front (trois ont une probabilité inférieure à 0,2), on perd de l'hypervolume que ce bilan ne compte pas. C'est ce qu'on observe sur le problème V2.

\textbf{Ex.~\ref{ex:clf-seuil}.} $\ln(0{,}1/0{,}9)=-2{,}20$ ; $\ln(0{,}2/0{,}8)=-1{,}39$ ; $\ln1=0$. Un log-odds de $-1{,}8$ n'est accepté qu'au seuil 0,1. Sa probabilité de succès est $1/(1+e^{1{,}8})\approx0{,}14$.

\section*{Chapitre~\ref{ch:hors-ligne} — Sans XFOIL}
\textbf{Ex.~\ref{ex:hl-exact}.} Le QUBO a $2^{11}=2\,048$ chaînes (10 bits de conception, 1 auxiliaire) et chaque résolution exacte en garde les 600 meilleures, soit près de 30~\%. L'optimum contraint, qui a un bon objectif, figure dans cette liste pour tous les $\lambda$ visités, et la vérification le retient puisqu'elle trie selon l'objectif seul. À 30 bits, il faudrait garder une fraction comparable de $10^9$ chaînes : c'est l'énumération, impossible. Le recuit ne rend que quelques dizaines de lectures distinctes, et ce sont les $\lambda$ qui décident où elles tombent.

\textbf{Ex.~\ref{ex:hl-gain}.} (a) La boucle ne simule que la conception ; dans le QUBO à 14 qubits, la même conception occupe 16 chaînes, dont une seule de basse énergie. Comparer des chaînes complètes avec des références $1/2^{14}$ et $1/2^{10}$ fausserait la comparaison. (b) Un gain de 1 veut dire que QAOA ne tire pas la conception optimale plus souvent qu'un tirage uniforme. Il peut faire mieux ailleurs (80~\% d'échantillons faisables, contre bien moins au hasard), mais il n'aide pas à trouver l'optimum. (c) Le cours IA et QUBO mesurait la probabilité de l'état fondamental du QUBO complet, chaîne entière, rapportée à $1/2^{14}$ ; cet état n'est pas toujours l'optimum contraint (petites violations rentables), et la normalisation diffère.

\section*{Chapitre~\ref{ch:cqm} — Interface}
\textbf{Ex.~\ref{ex:cqm-ecrire}.}
\begin{lstlisting}[language={}]
Minimize
 obj: b0 - 2 b1 + [ 6 b0 * b1 ] / 2
Subject To
 c1: b0 + b1 >= 1
 c2: - 0.5 b0 + [ b0 * b1 ] >= -0.2
Binary
 b0 b1
End
\end{lstlisting}
Énumération : $00$ viole c1 ; $10$ : objectif 1, c2 vaut $-0{,}5<-0{,}2$, infaisable ; $01$ : objectif $-2$, faisable ; $11$ : objectif $2$, faisable. L'optimum est $01$. Le crochet de l'objectif est divisé par 2 : on y écrit $6$ pour un couplage de 3.


\end{document}
